\documentclass{llncs}

\usepackage{amsmath}
\usepackage{amssymb}
%\usepackage{theorem}
\usepackage{graphicx}
\usepackage{mathpartir}
\usepackage{fancyvrb}
\usepackage{url}
\usepackage{ifpdf}

\ifpdf
\usepackage{pdfsync}
\fi

%%%% MACROS %%%%%

% Private comments, should be left in the source when the paper is finished.
% Use this version for visible comments:
\newcommand{\comment}[1]{\paragraph{Comment:} #1}
% For invisible comments:
%\newcommand{\comment}[1]{}

% Internal notes, used during writing the paper. The first argument is
% the author of the note. When the paper is finished, these should
% all be gone. Delete an internal note as soon as you resolve it.
\newcommand{\internal}[2]{\par\fbox{\parbox{0.9\textwidth}{{Note by #1:} #2}}\par}

% Types and terms in ocaml
\newcommand{\ctype}{\mathtt{type}\;}
\newcommand{\cint}{\mathtt{int}}
\newcommand{\clist}[1]{#1\;\mathtt{list}}
\newcommand{\ctrue}{\mathtt{true}}
\newcommand{\cwhile}[2]{\mathtt{while}\;#1\;\mathtt{do}\;#2\;\mathtt{done}}
\newcommand{\cstring}[1]{\mathtt{"#1"}}
\newcommand{\cprint}[1]{\mathtt{print\_string}\;#1}
\newcommand{\cfun}[2]{\mathtt{fun}\;#1 \to #2}
\newcommand{\clet}[2]{\mathtt{let}\;#1 = #2}
\newcommand{\cletrec}[2]{\mathtt{let}\;\mathtt{rec}\;#1 = #2}
\newcommand{\cfst}[1]{\mathtt{fst}\;#1}
\newcommand{\csnd}[1]{\mathtt{snd}\;#1}
\newcommand{\cmod}{\mathbin{\mathtt{mod}}}

% Realizaiblity
\newcommand{\per}{\approx}
\newcommand{\perty}[1]{\per_{\mathtt{#1}}}
\newcommand{\rz}{\Vdash}
\newcommand{\Type}{\mathsf{Type}}
\newcommand{\values}[1]{[\![#1]\!]}
\newcommand{\typeOf}[1]{|#1|}
\newcommand{\setOf}[1]{\langle #1 \rangle}
\newcommand{\support}[1]{\|#1\|}

% Sets
\newcommand{\set}[1]{\{#1\}}
\newcommand{\such}{\mid}
\newcommand{\pair}[1]{\langle #1 \rangle}
\newcommand{\epito}{\twoheadrightarrow}
\newcommand{\family}[2]{\set{#1}_{#2}}
\newcommand{\depsum}[2]{{\Sigma_{#1} #2}}
\newcommand{\depprod}[2]{{\Pi_{#1} #2}}

% Common sets
\newcommand{\NN}{\mathbb{N}}
\newcommand{\ZZ}{\mathbb{Z}}
\newcommand{\RR}{\mathbb{R}}

% Theorem-like environments
\iffalse
{
  \theorembodyfont{\itshape}
  \newtheorem{theorem}{Theorem}[section]
  \newtheorem{lemma}[theorem]{Lemma}
  \newtheorem{proposition}[theorem]{Proposition}
  \newtheorem{corollary}[theorem]{Corollary}
}
{
  \theorembodyfont{\rmfamily}
  \newtheorem{definition}[theorem]{Definition}
  \newtheorem{example}[theorem]{Example}
}
\fi

% Quantifiers. I am really used to these...
\newcommand{\all}[3]{\forall\, #1 \,{\in}\, #2\,.\left(#3\right)}
\newcommand{\some}[3]{\exists\, #1 \,{\in}\, #2\,.\left(#3\right)}
\newcommand{\exactlyone}[3]{\exists!\, #1 \,{\in}\, #2\,.\left(#3\right)}
\newcommand{\xall}[3]{\forall\, #1 \,{\in}\, #2\,.\,#3}
\newcommand{\xsome}[3]{\exists\, #1 \,{\in}\, #2\,.\,#3}
\newcommand{\xexactlyone}[3]{\exists!\, #1 \,{\in}\, #2\,.\,#3}


%%%%%%%%%%%%%%%%%%%%%%%%%%%%%%%%%%%%
%% Macros for input & output language

% ASCII syntax
\iffalse
\newcommand{\ip}{p}
\newcommand{\ix}{x}
\newcommand{\is}{s}
\newcommand{\ie}{e}
\newcommand{\iP}{P}
\newcommand{\iS}{S}
\newcommand{\iprop}{\mathtt{prop}}
\newcommand{\itrue}{\mathtt{true}}
\newcommand{\ifalse}{\mathtt{false}}
\newcommand{\inot}[1]{\mathtt{not}\,#1}
\newcommand{\iimply}[2]{#1\,\mbox{\texttt{->}}\,#2}
\newcommand{\iiff}[2]{#1\,\mbox{\texttt{<->}}\,#2}
\newcommand{\iequal}[2]{#1\,\mathtt{=}\,#2}
\newcommand{\iands}[2]{#1\,\mathtt{/\backslash}\, #2}
\newcommand{\iors}[2]{#1\,\mathtt{\backslash/}\,#2}
\newcommand{\iforall}[3]{\texttt{forall}\,#1\,\texttt{:}\,#2\texttt{,}\ #3}
\newcommand{\iexists}[3]{\texttt{exists}\,#1\,\texttt{:}\,#2\texttt{,}\ #3}
\newcommand{\iunique}[3]{\texttt{unique}\,#1\,\texttt{:}\,#2\texttt{,}\ #3}
\newcommand{\ilambda}[3]{\texttt{lambda}\,#1\,\texttt{:}\,#2\texttt{,}\ #3}
\newcommand{\iapps}[3]{#1\ #2\,\cdots\,#3}

\newcommand{\ituples}[2]{\texttt{(}#1\texttt{,}\ldots\texttt{,}#2\texttt{)}}
\newcommand{\iproj}[2]{#1\texttt{.}#2}
\newcommand{\isubin}[2]{#1\,\texttt{:>}\,#2}
\newcommand{\isubout}[2]{#1\,\texttt{:<}\,#2}
\newcommand{\ithe}[3]{\texttt{the}\,#1\,\texttt{:}\,#2\texttt{,}\, #3}
\newcommand{\ilet}[3]{\texttt{let}\,#1\,\texttt{=}\,#2\,\texttt{in}\, #3}
\newcommand{\iquot}[2]{#1\,\texttt{\%}\,#2}
\newcommand{\iletquot}[4]{\texttt{let}\ #1\texttt{\%}\,#2\,\texttt{=}\,#3\ \texttt{in}\ #4}
\newcommand{\irz}[1]{\texttt{rz}\,#1}
\newcommand{\iletrz}[3]{\texttt{let}\ \texttt{rz}\,#1\,\texttt{=}\,#2\ \texttt{in}\ #3}

\newcommand{\iprods}[2]{#1\,\texttt{*}\cdots\texttt{*}\,#2}
\newcommand{\isums}[2]{#1\,\texttt{+}\cdots\texttt{+}\,#2}

\fi

% Spiffed up syntax using mathematical symbols
\newcommand{\ip}{\varphi}
\newcommand{\ipp}{\rho}
\newcommand{\ix}{x}
\newcommand{\is}{s}
\newcommand{\ie}{e}
\newcommand{\iP}{p}
\newcommand{\iS}{\alpha}
\newcommand{\itrue}{\top}
\newcommand{\ifalse}{\bot}
\newcommand{\inot}[1]{\lnot #1}
\newcommand{\iimply}[2]{#1\Rightarrow #2}
\newcommand{\iiff}[2]{#1 \Leftrightarrow #2}
\newcommand{\iequal}[2]{#1 = #2}
\newcommand{\iand}[2]{#1\land #2}
\newcommand{\ior}[2]{#1\lor #2}
\newcommand{\iforall}[3]{\forall #1{:}#2.\ #3}
\newcommand{\iexists}[3]{\exists #1{:}#2.\ #3}
\newcommand{\iunique}[3]{\exists! #1{:}#2.\ #3}
\newcommand{\ilambda}[3]{\lambda #1{:}#2.\ #3}
\newcommand{\iapp}[2]{#1\ #2}

\newcommand{\ituple}[2]{(#1,\ldots,#2)}
\newcommand{\iproj}[2]{\mathop{\pi_{#2}} #1}
\newcommand{\isubin}[2]{#1\,{{:}{>}}\,#2}
\newcommand{\isubout}[2]{#1\,{{:}{<}}\,#2}
\newcommand{\ithe}[3]{\iota\,#1{:}#2{.}\, #3}
\newcommand{\ilet}[3]{\textsf{let}\,#1{=}#2\,\textsf{in}\, #3}
\newcommand{\ieclass}[2]{[#1]_{#2}}
\newcommand{\ileteclass}[4]{\textsf{let}\ [#1]_{#2}\,{=}#3\ \textsf{in}\ #4}
\newcommand{\irz}[1]{\textsf{rz}\,#1}
\newcommand{\iletrz}[3]{\textsf{let}\ \textsf{rz}\,#1\,{=}\,#2\ \textsf{in}\ #3}
\newcommand{\iinj}[2]{`{#1}\,#2}
\newcommand{\imatches}[7]{\textsf{match}\,#1\,\textsf{with}\,\iinj{#2}{#3}\Rightarrow{#4}\,|\cdots|\iinj{#5}{#6}\Rightarrow{#7}}
\newcommand{\imatch}[7]{\textsf{match}\,#1\,\textsf{with}\,\iinj{#2}{#3}\Rightarrow{#4}\,|\,\iinj{#5}{#6}\Rightarrow{#7}}
\newcommand{\iselect}[2]{#1{.}#2}

\newcommand{\itag}{\mathit{tag}}
\newcommand{\iprod}[3]{#1{:}#2\times#3}
\newcommand{\isum}[4]{#1{:}#2+#3{:}#4}
\newcommand{\iquot}[2]{#1{\bigm/} #2}
\newcommand{\isubset}[3]{\{#1{:}#2\ |\ #3\}}
\newcommand{\isunit}{1}
\newcommand{\isvoid}{0}

% I removed periods because I use these in the middle of text.
\newcommand{\iDefinition}[2]{\mathsf{Definition}\ #1 \mathbin{{:}{=}} #2}
\newcommand{\iParameter}[2]{\mathsf{Parameter}\ #1 : #2}
\newcommand{\iAxiom}[2]{\mathsf{Axiom}\ #1 : #2}
\newcommand{\ipt}{\tau}
\newcommand{\ik}{\kappa}
\newcommand{\iarrow}[2]{#1{\to}#2}
\newcommand{\idarrow}[3]{(#1{:}#2){\to}#3}
\newcommand{\ite}{\theta}
\newcommand{\ith}{\Theta}
\newcommand{\iTH}{T}
\newcommand{\ithy}[1]{\mathsf{thy}\ #1\ \mathsf{end}}
\newcommand{\im}{M}
\newcommand{\iM}{m}
\newcommand{\il}{l}
\newcommand{\iSet}{\mathsf{Set}}
\newcommand{\iStable}{\mathsf{Stable}}
\newcommand{\iProp}{\mathsf{Prop}}
\newcommand{\iEquiv}[1]{\mathsf{Equiv}(#1)}
\newcommand{\ibar}{\ |\ }

%%%%%%%%%%%%%%%%%%%%%%%%%%%%%%%%%%%%%%%%%%%%%%%%%%%%%%%%%%%%%%%%%%%%%%
%  Output-specific syntax (did not bother to produce ASCII version yet)
\newcommand{\obar}{\ |\ }
%%%% types
\newcommand{\oty}{\tau} % type
\newcommand{\opty}{\alpha} % polymorphic type
\newcommand{\oTY}{T} % type name
\newcommand{\oprod}[2]{#1 \times #2} % product
\newcommand{\oarrow}[2]{#1 \to #2} % arrow
\newcommand{\ounit}{\mathsf{unit}} % unit type
\newcommand{\oselect}[2]{#1{.}#2}  % projection from a module
\newcommand{\osumty}[4]{[`#1\,\mathsf{of}\,#2 \mid \cdots \mid `#3\,\mathsf{of}\,#4]}
\newcommand{\osumtyx}[4]{[`#1\,\mathsf{of}\,#2 \mid `#3\,\mathsf{of}\,#4]}
\newcommand{\ol}{l} % label
%%%% terms
\renewcommand{\oe}{e} % expression
\newcommand{\ox}{x} % variable
\newcommand{\olambda}[3]{\mathsf{fun}\,#1{:}#2 \to #3} % abstraction
\newcommand{\oapp}[2]{#1\,#2} % application
\newcommand{\otuple}[2]{(#1,\ldots,#2)} % tuple
\newcommand{\oproj}[2]{\mathsf{p}_{#2}\,#1} % projection
\newcommand{\oinj}[2]{`#1\,#2} % sum injection
\newcommand{\omatch}[7]{\mathsf{match}\,#1\,\mathsf{with}\,\oinj{#2}{#3}\to{#4}\obar\oinj{#5}{#6}\to{#7}}
\newcommand{\omatches}[7]{\mathsf{match}\,#1\,\mathsf{with}\,\oinj{#2}{#3}\to{#4}\,|\cdots|\oinj{#5}{#6}\to{#7}}
\newcommand{\ooblig}[4]{\mathsf{assure}\,#1{:}#2.\,#3\,\mathsf{in}\,#4}
\newcommand{\oobligx}[2]{\mathsf{assure}\,#1\,\mathsf{in}\,#2}
\newcommand{\olet}[3]{\mathsf{let}\,#1{=}#2\,\mathsf{in}\,#3}
%%%% propositions
\newcommand{\op}{p} % proposition
\newcommand{\oP}{P} % atomic proposition
\newcommand{\otrue}{\top} % true
\newcommand{\ofalse}{\bot} % false
\newcommand{\onot}[1]{\lnot #1} % not
\newcommand{\oand}[2]{#1 \land #2} % and
\newcommand{\oimply}[2]{#1 \Rightarrow #2} % imply
\newcommand{\oiff}[2]{#1 \Leftrightarrow #2} % iff
\newcommand{\ototal}[2]{#1{:}\support{#2}} % support membership
\newcommand{\oper}[3]{#2\,{\approx_{#1}}\,#3} % per
\newcommand{\oequal}[2]{#1{=}#2} % observational equality
\newcommand{\oforall}[3]{\forall #1{:}#2.\,#3}
\newcommand{\oforallt}[3]{\forall #1{:}\support{#2}.\,#3}
%%% Modest sets?
\newcommand{\oS}{S}
\newcommand{\os}{s}
%%% Modules
\newcommand{\om}{m} % module
\newcommand{\oM}{M} % module name
%%% Proposition kinds
\newcommand{\opt}{\Pi}
\newcommand{\oProp}{\mathsf{bool}}
%%% Module type elements
\newcommand{\ote}{\theta}
\newcommand{\ovalspec}[2]{\mathsf{val}\;#1:#2}
\newcommand{\otyspec}[1]{\mathsf{type}\;#1}
\newcommand{\otydef}[2]{\mathsf{type}\;#1=#2}
\newcommand{\omodulespec}[2]{\mathsf{module}\;#1:#2}
\newcommand{\osignatdef}[2]{\mathsf{module}\;\mathsf{type}\,#1=#2}
\newcommand{\opropspec}[2]{\mathsf{predicate}\; #1:#2}
\newcommand{\oassertion}[2]{\mathsf{assertion}\; #1 : #2}
\newcommand{\oA}{A} % Assertion name
%%% Module type
\newcommand{\omt}{\Sigma}
\newcommand{\oMT}{S}
\newcommand{\osig}[1]{\mathsf{sig}\,#1\,\mathsf{end}}
\newcommand{\ofunctor}[3]{\mathsf{functor}\, (#1{:}#2) \to #3}

%%% Local Variables: 
%%% mode: latex
%%% TeX-master: "cie"
%%% End: 


%%%%%%%%%%%%%%%%%%%%%%%%%%%%%%%%%%%%%%%%%%%%%%%%%%
%
% Slightly better macros for displaying source code.
% We can fiddle with the style later, I just stole this.
% (Use number=left to get source line numbers.)

% To include a source file
\newcommand{\sourcefile}[1]{%
\VerbatimInput[xleftmargin=1.5em,fontsize=\small]{#1}}

% To include lots of source
\newcommand{\sourcefilex}[1]{%
\VerbatimInput[xleftmargin=1.5em,fontsize=\scriptsize]{#1}}

% To show source
\DefineVerbatimEnvironment{source}{Verbatim}{xleftmargin=1.5em,fontsize=\small}

%%%%%%%%%%%%%%%%%%%%%%%%%%%%%%%%%%%%%%%%%%%%%%%%%%

\begin{document}
\title{RZ: a Tool for Bringing\\
  Constructive and Computable Mathematics\\
  Closer to Programming Practice}
\author{
  Andrej Bauer\inst{1} \and
  Christopher A. Stone\inst{2}}

\institute{
  Faculty of Mathematics and Physics, University of Ljubljana, Slovenia\\
  \email{Andrej.Bauer@fmf.uni-lj.si}
  \and
  Computer Science Department, 
  Harvey Mudd College, USA\\
  \email{stone@cs.hmc.edu}
}
\maketitle

\begin{abstract}
  We show that realizability is not only a fundamental tool in logic
  and computability, but also has direct application to the design and
  implementation of programs: it can produce a description of the data
  structure (a code interface) directly corresponding to a
  mathematical theory.

  Our claims are supported by an implementation of a tool, called RZ,
  which is to serve as a bridge between the world of constructive
  mathematics and the programming world. By using the realizability
  interpretation of constructive mathematics, RZ automatically
  translates specifications in constructive logic into standard
  interface code in a programming language, currently Objective Caml.

  RZ was specifically designed as a lightweight system, despite
  supporting a very rich input language that allows descriptions of
  complex mathematical structures. RZ does not extract code from
  proofs, but allows the programmer to implement the specifications in
  any fashion, ranging from hand-writen code to code extracted from
  proofs by other tools.
\end{abstract}

\section{Introduction}
\label{sec:introduction}

Given the constants, relations, and axioms of a mathematical structure,
expressed in constructive logic, what must a programmer do in order to
implement that mathematical structure?

For some familiar cases, the answer is obvious. Code implementing a
group must have a type to represent the group elements, a constant, a
unary operator, and a binary operator, each satisfying appropriate
axioms that guarantee the constant is a unit of the binary operator,
the unary operator computes inverses, and that the binary operator is
associative.

But for more interesting structures, especially those arising in
constructive analysis, the answer is much less obvious. Significant
research goes into finding suitable representations~\cite{???}. What
do you need to implement the real numbers (a complete, ordered field)?
What about a complete metric space? Or a [cs: I'm hoping Andrej will
fix this part up and add more interesting examples examples.] �`
\bigskip

The theory of realizability provides guidance in this task.
Realizability is a fundamental theoretical tool for the study of
logic, but it also has direct application to the design and
implementation of programs. It can be used to find a description of
the data structure directly corresponding to a mathematical
specification.

Extracting code interfaces by hand from mathematical axioms quickly
grows tedious, especially because different but logically equivalent
sets of axioms correspond to different, although interdefinable,
interfaces for code. One might want to see several of these, since
some interfaces can be more useful than others in practice.

And few programmers --- even those with solid backgrounds in
mathematics and logic --- are familiar with constructive logic or
realizability. Programmers are more familiar with language constructs
describing interfaces: C++ header files, ML signatures, Java
interfaces, and the like.

\bigskip

We have therefore implemented a system, called RZ, which serves as a
bridge between the logical world and the programming world. It
automatically translates specifications in constructive logic into
actual interface code in a programming language (currently Objective
Caml~\cite{ocaml}, but we are considering other languages as well). The
constructive part of the original specification turns into interface
code listing types and values that must exist in an implementation.

The remainder of the specification is maintained as assertions about
these types and values. Because these assertions have no computational
content, they can be interpreted constructively or classically equally
well, and in particular make equal sense to programmers (and
mathematicians) who are used to working in classical logic.

\bigskip

RZ was specifically designed as a lightweight system. Although the
realizability translation can be extended to transforming complete
proofs into complete code~\cite{komagata+:tr95}, we have not
implemented this. Other systems, including Coq~\cite{coq}, already
perform this task well.

But Coq works best when the entire task (from specification to code
generation) is performed within the same system. In contrast, the
interfaces generated by RZ can be implemented in any fashion, as long
as the assertions are satisfied. Code can be written by hand ---
allowing imperative, concurrent, and other arbitrary language
features, not just a ``functional'' subset. Or, the system could
generate a Coq interface as an \emph{input}, where the distinction
between computational (Set) and non-computational (Prop) is
automatically determined, and a corresponding implementation can be
provided through logical techniques.

The RZ system translates this to the language understood by many
programmers (typed interfaces with assertions in classical logic)

\bigskip

In this paper we show how RZ makes practical use of realizability and
other theoretical techniques, and provide some interesting motivational
examples.


%%% Local Variables: 
%%% mode: latex
%%% TeX-master: "cie"
%%% End: 

\section{Typed realizability}
\label{sec:typed-realizability}

RZ is based on \emph{typed realizability} by John
Longley~\cite{Longley99}.   This variant of realizability corresponds most
directly to programmers' intuition about implementations.

We motivate and explain typed realizability and its relationship to
real-world programming by way of example. Suppose we are asked to
design a data structure for the set $\mathcal{G}$ of all finite
simple\footnote{At most one arrow between any two vertices.}
directed graphs with vertices labeled by distinct integers. An exemplar
directed graph~$G$ is shown in Figure~\ref{fig:digraph}.
%
\begin{figure}
  \centering
  \includegraphics[width=0.4\textwidth]{digraph}
  \caption{A finite directed graph $G$}
  \label{fig:digraph}
\end{figure}
%
A common representation is \textbf{a} pair of lists $(\ell_V, \ell_A)$, where
$\ell_V$ is the list of vertex labels and $\ell_A$ is the \emph{adjacency list} 
representing the arrows by pairing the labels of each source and target.
In our example, $\ell_V = [1; 2;
3; 4]$ and $\ell_A = [(1,2); (2,2); (2,3); (3,2); (3;1)]$. Thus we
define the datatype of graphs to be\footnote{We use OCaml notation in
  which $\clist{t}$ classifies finite lists of elements of type~$t$, and
  $t_1 * t_2$ classifies pairs containing a value of type $t_1$ and and 
  value of type $t_2$.}
%
\begin{equation*}
  \ctype \mathtt{graph} = \clist{\cint} * \clist{(\cint * \cint)}
\end{equation*}
%
However, this is not a complete description of the representation, as
there are invariants and conditions which are not expressed directly
in this type, such as:
%
\begin{enumerate}
\item The order in which vertices and arrows are listed is not
  important; for example, $[1;2;3;4]$ and $[4;1;2;3]$ represent the same vertices.
\item Each vertex and arrow must be listed exactly once.
\item The source and target of each arrow must appear in the list of vertices.
\end{enumerate}
%
An implementation of the mathematical set~$\mathcal{G}$ must determine
not only what the underlying datatype $\mathtt{graph}$ is, but also
which values of type $\mathtt{graph}$ represent which elements
of~$\mathcal{G}$. As we shall see next, all of this can be expressed
either as a \emph{realizability relation} or a \emph{partial
  equivalence relation (per)}.

\subsection{Modest sets and pers}
\label{sec:modest-sets-pers}

We now give the formal definition of typed realizability, as it
applies to OCaml. Other general-purpose programming languages could be
used instead, as long as they provide the usual ground types, product
and function types.\footnote{It is also convenient to work with a
language that supports sum types, as this allows a more natural
representation of disjoint unions.}

Let $\Type$ be the collection of all (non-parametric) OCaml types. To
each type $t \in \Type$ we assign the set $\values{t}$ of values of
type~$t$ which behave \emph{functionally} in the sense
of~\cite{longley99when}. Such values are represented by terminating
expressions which do not raise exceptions or return different results
on different invocations, although they \emph{may} use exceptions,
store, and other computational effects, provided they behave from outside
as if they did not. A useful example of a functional value using store is
presented in Section~\ref{sec:we-show-modulus-of-continuity-example}.
Note also that a functional value with a function type may diverge as
soon as it is applied, e.g., if we define
$\cletrec{f\;(x:\cint):\cint}{f\;x}$ then $f \in \values{\cint \to
  \cint}$. The collection $\Type$ with the assignment of functional
values $|t|$ to each $t \in \Type$ forms a \emph{typed partial
  combinatory algebra (TPCA)}, which provides a theoretical basis for
the definition of a realizability model that suits our needs.

Going back to our example, we see that an implementation of 
directed graphs $\mathcal{G}$ specifies a datatype
$\typeOf{\mathcal{G}} = \mathtt{graph}$ together with a
\emph{realizability relation} $\rz_{\mathcal{G}}$ between
$\mathcal{G}$ and $\values{\mathtt{graph}}$. The meaning of $(\ell_V,
\ell_A) \rz_\mathcal{G} G$ is ``OCaml value $(\ell_V, \ell_A)$ 
represents/realizes/implements graph $G$''. There are two natural conditions
that $\rz_\mathcal{G}$ ought to satisfy: (1) for every $G \in
\mathcal{G}$ there should be at least one realizer $(\ell_V, \ell_A)$
representing it, and (2) if $(\ell_V, \ell_A)$ represents both $G$ and
$G'$ then $G = G'$. The latter condition is called \emph{modesty} and
is not strictly necessary for the development of the theory, though it
is something that programmers would naturally expect to hold. If
$(\ell_V, \ell_A)$ and $(\ell'_V, \ell'_A)$ represent the same graph 
(e.g., because $\ell_V$ is a permutation of $\ell'_V$, and similarly for $\ell_A$ and $\ell'_A$)
we say that they are \emph{equivalent} and write $(\ell_V, \ell_A)
\per_\mathcal{G} (\ell'_V, \ell'_A)$. The relation $\per_\mathcal{G}$
is a \emph{partial} equivalence relation (symmetric and transitive,
but not reflexive) because not every $(\ell_V, \ell_A) \in
\values{\mathtt{graph}}$ represents a graph.

\bigskip

A general definition is in order. A \emph{modest set} is a triple $A =
(\setOf{A}, \typeOf{A}, {\rz_A})$ where $\setOf{A}$ is the
\emph{underlying set}, $\typeOf{A} \in \Type$ is the \emph{underlying
  type}, and $\rz_A$ is a \emph{realizability relation} between
$\values{\typeOf{A}}$ and $\setOf{A}$, satisfying
% 
\begin{enumerate}
\item \emph{totality:} for every $x \in \setOf{A}$ there is $v \in
  \values{\typeOf{A}}$ such that $v \rz_A x$, and
\item \emph{modesty:} if $u \rz_A x$ and $u \rz_A y$ then $x = y$.
\end{enumerate}
%
The \emph{support} of $A$ is the set $\support{A} = \set{v \in
  \values{\typeOf{A}} \such \xsome{x}{\setOf{A}}{v \rz_A x}}$ of those
values which realize something. We define the relation $\per_A$ on
$\values{\typeOf{A}}$ by
%
\begin{equation*}
  u \per_A v
  \iff
  \some{x}{\setOf{A}}{u \rz_A x \land v \rz_A x} \;.
\end{equation*}
%
From totality and modesty of $\rz_A$ it follows that $\per_A$ is a per,
i.e., symmetric and transitive. Observe that $\support{A} = \set{v \in
  \values{\typeOf{A}} \such v \per_A v}$, whence $\per_A$
restricted to $\support{A}$ is an equivalence relation. In fact, we
may recover a modest set up to isomorphism from $\typeOf{A}$ and
$\per_A$ by taking $\setOf{A}$ to be the set of equivalence classes of
$\per_A$, and $v \rz_A x$ to mean $v \in x$.

The two views of implementations, as modest sets $(\setOf{A},
\typeOf{A}, {\rz_A})$, and as pers $(\typeOf{A}, {\per_A})$, are
equivalent.\footnote{And there is a third one, as a partial surjection
  $\delta_A : {\subseteq}\values{\typeOf{A}} \epito \setOf{A}$, with
  $\delta_A(v) = x$ when $v \rz_A x$. This is how realizability is
  presented in Type Two Effectivity~\cite{Wei00}.} In RZ we use pers
because they refer only to types and values, as opposed to arbitrary
sets. Nevertheless, it is useful to understand how modest sets and
pers arise from natural considerations about programming practice.

Modest sets form a category whose objects are modest sets and
morphisms are the realized functions. A \emph{realized function} $f :
A \to B$ is a function $f : \setOf{A} \to \setOf{B}$ for which that there
exists $v \in \values{\typeOf{A} \to \typeOf{B}}$ such that, for all
$x \in \setOf{A}$ and $u \in \typeOf{A}$,
%
\begin{equation}
  \label{eq:rz-function-space}
  u \rz_A x \implies v\;u \rz_B f(x) \;.
\end{equation}
%
This condition is just a mathematical expression of the usual idea
that~$v$ is an implementation of~$f$, since~$v$ does to realizers
what~$f$ does to the elements they represent.

The equivalent category of pers has as objects pairs $A = (\typeOf{A},
{\per_A})$ where $\typeOf{A} \in \Type$ and $\per_A$ is a per on
$\values{\typeOf{A}}$. A morphism $A \to B$ is represented by $v
\in \values{\typeOf{A} \to \typeOf{B}}$ such that, for all $u, u'
\in \support{A}$,
%
\begin{equation}
  \label{eq:per-exponential}
  u \per_A u' \implies v\;u \per_B v\;u' \;.
\end{equation}
%
Values $v$ and $v'$ which both satisfy~\eqref{eq:per-exponential}
represent the same morphism if, for all $u, u' \in \support{A}$,
%
\begin{equation*}
  u \per_A u' \implies v\;u \per_B v'\;u' \;.
\end{equation*}

The category of modest sets has a very rich structure. For example, we
may form a cartesian product $A \times B$ of modest sets $A$ and $B$
by
%
\begin{align*}
  \setOf{A \times B} &= \setOf{A} \times \setOf{B},\\
  \typeOf{A \times B} &= \typeOf{A} * \typeOf{B},\\
  p \rz_{A \times B} (x,y) &\iff
  \cfst{p} \rz_A x \land \csnd{p} \rz_B y.
\end{align*}
%
The projections $\pi_1 : A \times B \to A$ and $\pi_2 : A \times B \to
B$ are realized by $\mathtt{fst}$ and $\mathtt{snd}$, respectively.

The morphisms between modest sets~$A$ and~$B$ again form a modest set
$B^A$, also written as $A \to B$, called the \emph{exponential} of~$A$
and~$B$, with the underlying set
%
\begin{equation*}
  \setOf{B^A} =
  \set{f : \setOf{A} \to \setOf{B} \such \text{$f$ is a realized function}},
\end{equation*}
%
the underlying type
%
\begin{equation*}
  \typeOf{B^A} = \typeOf{A} \to \typeOf{B},
\end{equation*}
%
and the realizability relation $\rz_{B^A}$ defined
by~\eqref{eq:rz-function-space}. The evaluation map $e : B^A \times A
\to B$, $e(f,x) = f(x)$ is realized by OCaml application,
$\cfun{(u,v)}{u\;v}$. If a function $f : C \times A \to B$ is realized
by $v$, then its transpose $\tilde{f} : C \to B^A$, $\tilde{f}(z)(x) =
f(z,x)$, is realized by $\cfun{z\;x}{v \; (z,x)}$. This shows that the
category of modest sets is cartesian closed. In
Section~\ref{sec:transl-sets-terms} we review other canonical
constructions on modest sets.

\bigskip

As an example we consider the cyclic group on seven elements $(\ZZ_7,
0, {-}, {+})$. To implement the group, we must give a representation
of $\ZZ_7$ as a modest set~$Z = (\ZZ_7, \typeOf{Z}, {\rz_Z})$, and
provide realizers for the neutral element~$0$, negation~$-$, and
addition~$+$. We represent the elements of~$\ZZ_7 = \set{0, 1, \ldots,
  6}$ with values of type $\typeOf{Z} = \cint$. We allow any integer
$v \in \values{\cint}$ to represent the element $v
\mathbin{\mathrm{mod}} 7$,
%
\begin{equation*}
  v \rz_Z k \iff v \mathbin{\mathrm{mod}} 7 = k \;.
\end{equation*}
%
The neutral elements $0$ is represented by~$0$, while negation and
addition are respectively implemented by
%
\begin{align*}
  &\clet{\mathtt{neg}\; k}{-(k \cmod 7)} \\
  &\clet{\mathtt{add}\; (k, m)}{(k \cmod 7) + (m \cmod 7)}
\end{align*}
%
The expected $\clet{\mathtt{neg}\; k}{-k}$ and $\clet{\mathtt{add}\;
  k\; m}{k+m}$ do not work because of overflow at values that exceed
$2^{30}$ (in OCaml on a 32-bit machine). 

We could avoid the overflow worries by choosing a different, although
isomorphic, implementation in which the only valid realizers are $0$,
$1$, \ldots, $6$, and negation and addition are realized by
$\clet{\mathtt{neg}\; k}{7 - k}$ and $\clet{\mathtt{add}\; k\; m}{(k +
  m) \cmod 7}$, respectively.

This particular modest set~$Z$ also has \emph{decidable equality},
which means that the characteristic map $\mathrm{eq} : \ZZ_7 \times
\ZZ_7 \to \set{\mathtt{false}, \mathtt{true}}$ of equality is
realized, namely by\footnote{Incidentally, $\cfun{(k,m)}{(k - m) \cmod
    7 = 0}$ does not decide equality because of possible overflow in
  subtraction, whereas $\cfun{(k,m)}{k \cmod 7 = m \cmod 7}$ does not
  work because in OCaml $a \cmod b$ is negative when~$a$ is negative,
  e.g., $(-10) \cmod 7 = -3$ instead of the expected $4$.}
$\cfun{(k,m)}{(k \cmod 7 - m \cmod 7) \cmod 7 = 0}$.

However, not all modest sets have decidable
equality; there may not be a decision procedure for determining
whether two realizers represent the same element. For example, this
would happen if we implemented a semigroup with an undecidable word
problem~\cite{post47:_recur_unsol_probl_thue}.


\subsection{Interpretation of logic}
\label{sec:interpretation-logic}

In the realizability interpretation of logic, each formula~$\phi$ is
assigned a set of \emph{realizers} which can be thought of as
computations that witness the validity of~$\phi$. The situation is
somewhat similar, but not equivalent, to the propositions-as-types
translation of logic into type theory, where the proofs of a
proposition correspond to terms of the corresponding type. More
precisely, to each formula~$\phi$ we assign an underlying type
$\typeOf{\phi}$ of realizers, but unlike in the propositions-as-types
translation, not all terms of type $\typeOf{\phi}$ are necessarily
valid realizers for~$\phi$, and some terms which are realizers may not
correspond to any proofs, for example, if they denote partial
functions or use computational effects.

We write $t \rz \phi$ when $t \in \values{\typeOf{\phi}}$ is a
realizer for~$\phi$. The underlying types and the realizability
relation~$\rz$ are defined inductively on the structure of~$\phi$; an
outline is shown in Figure~\ref{fig:rz-logic}. We say that a
formula~$\phi$ is \emph{valid} if it has at least one realizer.
%
\begin{figure*}
  \textbf{Underlying types of realizers:}
  \begin{xalignat*}{2}
    \typeOf{\itrue} &= \mathtt{unit} &
    \typeOf{\ifalse} &= \mathtt{unit} \\
    \typeOf{\iequal{x}{y}} &= \mathtt{unit} &
    \typeOf{\iand{\phi}{\psi}} &= \oprod{\typeOf{\phi}}{\typeOf{\psi}} \\
    \typeOf{\iimply{\phi}{\psi}} &= \oarrow{\typeOf{\phi}}{\typeOf{\psi}} &
    \typeOf{\ior{\phi}{\psi}} &=
    \osumtyx{\mathtt{or}_1}{\typeOf{\phi_1}}{\mathtt{or}_2}{\typeOf{\phi_2}} \\
    \typeOf{\iforall{x}{A}{\phi}} &= \oarrow{\typeOf{A}}{\typeOf{\phi}} &
    \typeOf{\iexists{x}{A}{\phi}} &= \oprod{\typeOf{A}}{\typeOf{\phi}}
  \end{xalignat*}
  \textbf{Realizers:}
  \begin{align*}
    () \rz \itrue & \\
    () \rz \iequal{x}{y}
    &\quad\text{iff}\quad 
    x = y
    \\
    (t_1,t_2) \rz \iand{\phi}{\psi}
    &\quad\text{iff}\quad
    \text{$t_1 \rz \phi$ and $t_2 \rz \psi$}
    \\
    t \rz \iimply{\phi}{\psi}
    &\quad\text{iff}\quad
    \text{for all $u \in \typeOf{\phi}$, if $u \rz \phi$ then $t\,u
      \rz \psi$}
    \\
    \oinj{\mathtt{or}_1}{t} \rz \ior{\phi}{\psi}
    &\quad\text{iff}\quad
    \text{$t \rz \phi$}
    \\
    \oinj{\mathtt{or}_2}{t} \rz \ior{\phi}{\psi}
    &\quad\text{iff}\quad
    \text{$t \rz \psi$}
    \\
    t \rz \iforall{x}{A}{\phi(x)}
    &\quad\text{iff}\quad
    \text{for all $u \in \typeOf{A}$, if $u \rz_A x$ then $t\,u \rz \phi(x)$}
    \\
    (t_1, t_2) \rz \iexists{x}{A}{\phi(x)}
    &\quad\text{iff}\quad
    \text{$t_1 \rz_A x$ and $t_2 \rz \phi(x)$}
  \end{align*}
  \caption{Realizability interpretation of logic (outline)}
  \label{fig:rz-logic}
\end{figure*}

In classical mathematics, a predicate on a set~$X$ may be viewed as a
subset of~$X$ or a (possibly non-computable) function $X \to \oProp$,
where $\oProp = \set{\ofalse, \otrue}$ is the set of truth values.
Accordingly, since in realizability propositions are witnessed by
realizers, a predicate~$\phi$ on a modest set~$A$ may be viewed as a
subset of $\setOf{A} \times \values{\typeOf{\phi}}$, or a (possibly
non-computable) function $\setOf{A} \times \values{\typeOf{\phi}} \to
\set{\ofalse, \otrue}$. In terms of pers, which is what RZ uses, a
predicate~$\phi$ on a per~$A = (\typeOf{A}, {\per_A})$ is a (possibly
non-computable) function $\phi : \values{\typeOf{A}} \times
\values{\typeOf{\phi}} \to \oProp$ which is
%
\begin{itemize}
\item \emph{strict:} if $\phi(u,v)$ then $u \in \support{A}$, and
\item \emph{extensional:} if $\phi(u_1,v)$ and $u_1 \per_A u_2$ then
  $\phi(u_2,v)$.
\end{itemize}

We illustrate how the realizability interpretation extracts the
computational content of a proposition. To make an interesting
example, suppose we already implemented the structure of real
numbers~$\RR$ as a modest set~$R = (\RR, \mathtt{real}, {\rz_R})$, and
consider the statement that every cubic $x^3 + a x + b$ has a root,
%
\begin{equation}
  \label{eq:square-root}%
  \iforall{a}{R}{\iforall{b}{R}{\iexists{x}{R}{\iequal{x^3 + a x + b}{0}}}}.
\end{equation}
%
By computing according to Figure~\ref{fig:rz-logic}, we find out that
a realizer is a value~$r$ of type
$\oarrow{\mathtt{real}}{\oarrow{\mathtt{real}}{\oprod{\mathtt{real}}{\mathtt{unit}}}}$
such that, whenever $t$ realizes $a \in \RR$ and $u$ realizes $b \in
\RR$, then $\oapp{\oapp{r}{t}}{u} = (v, w)$ with $v$ realizing a real
number~$x$ such that $x^3 + a x + b = 0$, and $w$ is irrelevant. This
can be ``optimized'' to a realizer of type
$\oarrow{\mathtt{real}}{\oarrow{\mathtt{real}}{\mathtt{real}}}$ which
does not bother to compute~$w$. In essence then, a realizer~$r$
for~\eqref{eq:square-root} computes a root of the cubic equation. Note
that $r$ is \emph{not} extensional, i.e., different realizers~$t$
and~$u$ for the same~$a$ and~$b$ may result in different roots of the
cubic. To put this in another way, $r$ realizes a \emph{multi-valued}
function\footnote{The multi-valued nature of the realizer comes from
  the fact that it computes \emph{any one} of many values, not that it
  computes \emph{all} of the many values.} rather than a modest set
morphism. It is well known in computable mathematics that certain
operations, such as equation solving, are only computable if we allow
them to be multi-valued. They arise naturally in RZ as translations of
$\forall\exists$~statements.

There are propositions whose realizers are ``irrelevant'' or free of
computational content. For example, realizers for $\itrue$ and
equality have type $\ounit$. Another example is a negation
$\inot{\phi}$, which is defined to be the same as
$\iimply{\phi}{\ifalse}$, whose realizers have type
$\oarrow{\typeOf{\phi}}{\ounit}$. Such realizers do not compute
anything useful, and we may as well throw them away. Sometimes only a
part of a realizer is computationally irrelevant, as we saw in the
last example. Propositions which are free of any computational content
are characterized as the \emph{$\lnot\lnot$-stable propositions}. A
proposition~$\phi$ is said to be $\lnot\lnot$-stable, or just
\emph{stable} for short, when $\iimply{\inot{\inot{\phi}}}{\phi}$ is
valid. Any \emph{negative} proposition, i.e., one built from $\itrue$,
$\ifalse$, $=$, $\land$, $\Rightarrow$ and $\forall$ is stable, but
there may be other propositions which are stable and are not written
in the negative form.

It would be unproductive to bother the programmer with requirements
for useless code.  On input, RZ allows users to identify basic propositions
as computationally irrelevant.  The extracted
realizers also go through a \emph{thinning} phase, see Section~\ref{sec:translation}, which
optimizes away references to irrelevant realizers.

\internal{Chris}{It seemed out of place to introduce
the idea of subkinds and syntax for $\iStable$ and $\iProp$ here.}

\subsection{Uniform families of modest sets}
\label{sec:uniform-families}

Many structures are naturally viewed as families of sets, or sets
depending on parameters, or \emph{dependent types} as they are called
in type theory. For example, the $n$-dimensional Euclidean space
$\RR^n$ depends on the dimension $n \in \NN$, the Banach space
$\mathcal{C}([a,b])$ of uniformly continuous real functions on the
closed interval $[a,b]$ depends on $a, b \in \RR$ such that $a < b$,
etc. In general, a family of sets $\family{A_i}{i \in I}$ is an
assignment of a set $A_i$ to each $i \in I$ from an \emph{index
  set}~$I$.

In the category of modest sets the appropriate notion is that of a
\emph{uniform} family~$\family{A_i}{i \in I}$, which is an assignment
of a modest set $A_i = (\setOf{A_i}, \typeOf{A}, {\rz_{A_i}})$ to each
$i \in \setOf{I}$, where $I$ is an index modest set,
cf.~\cite{Jacobs,Birkedal}. The uniformity comes from the requirement
that all the~$A_i$'s share the same underlying type~$\typeOf{A_i} =
\typeOf{A}$. It is a desirable restriction from the implementation
point of view, because it removes dependencies at the level of types.
Note also that there is no dependency on the realizers, only on the
elements of the underlying set.

We may express uniform families in terms of pers, just like we did for
predicates. A uniform family $\family{A_i}{i \in I}$ of pers indexed
by a per~$I$ is given by an underlying type $\typeOf{A}$ and a family
of pers $(\per_{A_i})_{i \in \values{\typeOf{I}}}$ which is
% 
\begin{itemize}
\item \emph{strict:} if $u \per_{A_i} v$ then $i \in \support{I}$, and
\item \emph{extensional:} if $u \per_{A_i} v$ and $i \per_I j$ then $u
  \per_{A_j} v$.
\end{itemize}
%
A morphism $f : \family{A_i}{i \in I} \to \family{B_i}{i \in I}$ is a
family $\family{f_i : \setOf{A_i} \to \setOf{B_i}}{i \in I}$ of
uniformly realized maps, which means that there exists $u \in
\values{\typeOf{I} \to \typeOf{A} \to \typeOf{B}}$ such that,
%
\begin{equation*}
  v \rz_I i \land w \rz_{A_i} x \implies u\; v \; w \rz_{B_i} f_i(x)
\end{equation*}
%
for all $i \in I$, $v \in \values{\typeOf{I}}$, $x \in \setOf{A_i}$,
and $w \in \values{\typeOf{A}}$.

We may form the \emph{sum} $\depsum{i \in I}{A_i}$ of a uniform family
$\family{A_i}{i \in I}$ as
%
\begin{align*}
  \setOf{\depsum{i \in I}{A_i}} &=
  \set{\pair{i,x} \such i \in \setOf{I} \land x \in \setOf{A_i}}
  \\
  \typeOf{\depsum{i \in I}{A_i}} &=
  \typeOf{I} \times \typeOf{A}
  \\
  (u,v) \rz_{\depsum{i \in I}{A_i}} (i,x)
  &\iff
  u \rz_I i \land v \rz_{A_i} x
\end{align*}
%
and the \emph{product} $\depprod{i \in I}{A_i}$ as
%
\begin{align*}
  \setOf{\depprod{i \in I}{A_i}} &=
  \set{f : \setOf{I} \to {\textstyle \bigcup_{i \in \setOf{I}} \setOf{A_i}} \such
    \xall{i}{\setOf{i}}{f(i) \in \setOf{A_i}}}
  \\
  \typeOf{\depprod{i \in I}{A_i}} &=
  \typeOf{I} \to \typeOf{A}
  \\
  u \rz_{\depprod{i \in I}{A_i}} f
  &\iff
  \xall{i}{I}{\all{v}{\typeOf{I}}{
      v \rz_I i \implies
      u \; v \rz_{A_i} f(i)
    }}.
\end{align*}
%
These constructions allow us to interpret (extensional) dependent type
theory in the category of modest sets.

As an example of a uniform family we consider the cyclic group
$(\ZZ_n, 0, {-}, {+})$ of order~$n$. To keep things simple, we assume
that~$n$ ranges over natural numbers that can be represented by
type~$\cint$. More precisely, let $N$ be the modest set
%
\begin{align*}
  & \setOf{N} = \set{1, \ldots, 2^{30}} \\
  & \typeOf{N} = \cint \\
  & v \rz_N n \iff v = n.
\end{align*}
%
The uniform family $\family{Z_n}{n \in N}$ is like the above example
of the cyclic group of order~$7$, with $7$ replaced by~$n$:
%
\begin{align*}
  & \setOf{Z_n} = \ZZ_n = \set{0, 1, \ldots, n - 1}
  \\
  & \typeOf{Z_n} = \cint
  \\
  & v \rz_{Z_n} k \iff v \mathbin{\mathrm{mod}} n = k
  \\
  & \clet{\mathtt{neg} \; n \; k }{- (k \cmod n)}
  \\
  & \clet{\mathtt{add} \; n \; (k,m)}{k \cmod n + m \cmod n}.
\end{align*}
%
Another change is the additional parameter~$n$ accepted by the
realizers for negation and addition. The definitions make sense for
any $n \in \setOf{N}$, but for too large a value overflow might occur.
If we expect the implementation to be used with such values, it is
more prudent to replace the type~$\cint$ with an implementation of
arbitrary size integers.

%%% Local Variables: 
%%% mode: latex
%%% TeX-master: "cie"
%%% End: 

\section{Specifications as signatures with assertions}
\label{sec:spec-sign-assert}

\iflong
\begin{figure}[t]
  \centering
  \sourcefile{group.ml}
  \caption{The module type $\mathtt{Ab}$.}
  \label{fig:module-example}
\end{figure}
\fi % \iflong

\iflong
\begin{figure}
  \[
  \begin{array}{rl@{\qquad}l}
    \noalign{\textbf{Types}}
    \oty ::= 
    & \oTY  \obar \oselect{\oM}{\oTY} &\mbox{Type names}\\
    | & \ounit \obar \oprod{\oty_1}{\oty_2} &\mbox{Unit and cartesian product}\\
    | & \oarrow{\oty_1}{\oty_2} & \mbox{Function type}\\
    | & \osumty{\ol_1}{\oty_1}{\ol_n}{\oty_n} &\mbox{Disjoint sum}\\
    | & \opty & \mbox{Polymorphic types}\\[5pt]
    
    \noalign{\textbf{Terms}}	
    \oe ::=
    & \ox \obar \oselect{\oM}{\ox} &\mbox{Term names}\\
    | & \olambda{\ox}{\oty_1}{\oe} \obar 
    \oapp{\oe_1}{\oe_2} &\mbox{Functions and application}\\
    | & ()
    \obar \otuple{\oe_1}{\oe_n} 
    \obar \oproj{\oe}{n}&\mbox{Tuples and projection}\\
    | & \oinj{\ol}{\oe} 
    \obar (\omatches{\oe}{\ol_1}{\ox_1}{\oe_1}{\ol_n}{\ox_n}{\oe_n})&\mbox{Injection and projection from a sum}\\
    | & \ooblig{\ox}{\oty}{\op}{\oe} \obar \oobligx{\op}{\oe} &\mbox{Obligations}\\
    | & \olet{\ox}{\oe_1}{\oe_2}&\mbox{Local definitions}\\[5pt]
    

    \noalign{\textbf{Propositions (negative fragment)}}
    \op ::= 
    & \oP \obar \oselect{\oM}{\oP}&\mbox{Atomic
      proposition}\\
    | & \otrue  \obar \ofalse \obar \onot{\op_1} \obar \oand{\op_1}{\op_2} \obar 
    \oimply{\op_1}{\op_2} \obar \oiff{\op_1}{\op_2} & \mbox{Predicate logic}\\ 
    | & \opcases{\oe}{\ol_1}{\ox_1}{\op_1}{\ol_n}{\op_n}{\oe_n}&\mbox{Propositional case}\\
    | & \olambda{\ox}{\oty}{\op}  \obar \oapp{\op}{\oe} &
    \mbox{Propositional functions and application}\\
    | & \oper{\os}{\oe_1}{\oe_2} \obar \ototal{\oe}{\os} & \mbox{Pers and support} \\
    | & \oequal{\oe_1}{\oe_2} & \mbox{(Observational) term equality}\\
    | & \oforall{\ox}{\oty}{\op}  \obar 
    \oforallt{\ox}{\os}{\op} & \mbox{Term quantifiers}\\[5pt]
    
    \noalign{\textbf{Basic modest sets}}
    \os ::=
    & \oS \obar \oapp{\os}{\oe} & \\[5pt]

    \noalign{\textbf{Modules}}		
    \om ::= 
    & \oM  \obar \oselect{\om}{\oM}&\mbox{Model names}\\
    | & \oapp{\om_1}{\om_2}&\mbox{Application of parameterized model}\\[5pt]
    
    \noalign{\textbf{Proposition Kinds}}
    \opt ::=
    & \oProp & \mbox{Classifier for propositions}\\
    |& \oarrow{\oty}{\opt} & \mbox{Classifier for a predicate/relation}\\[5pt] 
    
    \noalign{\textbf{Specification elements}}
    \ote ::=
    & \ovalspec{\ox}{\oty}&\mbox{Value declaration}\\
    | & \otyspec{\oTY}&\mbox{Type declaration}\\
    | & \otydef{\oTY}{\oty}&\mbox{Type definition}\\
    | & \omodulespec{\om}{\omt}&\mbox{Module declaration}\\
    | & \osignatdef{\oMT}{\omt}&\mbox{Specification definition}\\
    | & \opropspec{\oP}{\opt}&\mbox{Proposition declaration}\\
    | & \oassertion{\oA}{\op}&\mbox{Assertion}\\[5pt]

    \noalign{\textbf{Specifications (module types with assertions)}}
    \omt ::= 
    & \oMT \obar \oselect{\om}{\oMT} & \mbox{Specification names}\\
    | & \osig{\ote_1 \ldots \ote_n} & \mbox{Specification elements}\\
    | & \ofunctor{\om}{\omt_1}{\omt_2} & \mbox{Parameterized specification}\\
%    | & \oapp{\omt}{\om} & \mbox{Specification application}
  \end{array}
  \]
  \label{fig:output}
  \caption{The syntax of specifications (Simplified)}
\end{figure}
\fi % \iflong

In programming we distinguish between \emph{implementation} and
\emph{specification} of a structure. In OCaml these two notions are
expressed with modules and module types, respectively.%
\iflong
\footnote{In
  object-oriented languages implementations and specifications are
  expressed with classes and interfaces, while in Haskell they
  correspond to modules and declarations.}
\else % \iflong
\ 
\fi % \iflong
A module defines types and
values, while a module type simply lists the types, type definitions,
and values provided by a module. For a complete specification, a
module type must also be annotated with \emph{assertions} which
specify the required properties of declared types and values.
\iflong
%
For example, if we look at the definition of module type $\mathtt{Ab}$
in Figure~\ref{fig:module-example}, we might guess 
that $\mathtt{Ab}$ is a signature for an Abelian group. 
%
However, $\mathtt{Ab}$ by itself requires an implementation satisfy
only the signature of an Abelian
group, but does not guarantee it behaves as an Abelian group.
A complete description would contain the
following further \emph{assertions}:
%
\begin{enumerate}
\item
  \label{enum:t-per}%
  there is a per $\perty{t}$ on $\values{\mathtt{t}}$,
\item
  \label{enum:zero-total}%
  $\mathtt{zero} \in \support{\mathtt{t}}$.
\item
  \label{enum:neg-total}%
  for all $u, v \in \values{\mathtt{t}}$, if $u \perty{t} v$ then
  $\mathtt{neg} \; u \perty{t} \mathtt{neg} \; v$,
\item
  \label{enum:add-total}%
  for all $u_1, u_2, v_1, v_2 \in \values{\mathtt{t}}$, if $u_1
  \perty{t} v_1$ and $u_2 \perty{t} v_2$ then $\mathtt{add} \; (u_1,
  u_2) \perty{t} \mathtt{add} \; (v_1, v_2)$,
\item 
  \label{enum:ab-group-axiom-1}%
  for all $u \in \support{\mathtt{t}}$, $\mathtt{add} \;
  (\mathtt{zero}, u) \perty{t} u$,
\item
  \label{enum:ab-group-axiom-2}%
  for all $u \in \support{\mathtt{t}}$, $\mathtt{add} \; (u,
  \mathtt{neg} \; u) \perty{t} \mathtt{zero}$,
\item
  \label{enum:ab-group-axiom-3}%
  for all $u, v, w \in \support{\mathtt{t}}$, $\mathtt{add} \;
  (\mathtt{add} \; (u, v), w) \perty{t} \mathtt{add} \; (u,
  \mathtt{add} \; (v, w))$,
\item
  \label{enum:ab-group-axiom-4}%
  for all $u, v \in \support{\mathtt{t}}$, $\mathtt{add} \; (u, v)
  \perty{t} \mathtt{add} (v, u)$.
\end{enumerate}
%
Assertions \ref{enum:zero-total}--\ref{enum:add-total} state that
$\mathtt{zero}$, $\mathtt{neg}$, and $\mathtt{add}$ realize a
constant, a unary, and a binary operation, respectively, while
assertions \ref{enum:ab-group-axiom-1}--\ref{enum:ab-group-axiom-4}
correspond to axioms for Abelian groups. 

The output of RZ consists of \emph{module specifications}, module types
plus assertions about their components.  More specifically, a specification
may contain value declarations, type declarations and definitions, module
declarations, specification definitions, proposition declarations, and
assertions.  The language of specifications is summarized in
Figure~\ref{fig:output}.

A specification can describe an OCaml structure (a collection of 
definitions for types and values) or an OCaml functor (a parameterized
module, i.e., a function mapping modules to modules).  The latter would
be appropriate, for example, when describing a uniform implementation of the real numbers
that works given any implementation of natural numbers.

Assertions within module specifications (which appear as code comments)
are expressed in the \emph{negative fragment} of
first-order logic, which contains constants for truth and falsehood,
negation, conjunction, implication, equivalence, and universal
quantification (but no disjunction or existential). 
%
\iflong
For convenience we
also introduce the propositional case
%
\begin{equation*}
  \opcases{\oe}{\ol_1}{\ox_1}{\op_1}{\ol_n}{\ox_n}{\op_n}
\end{equation*}
%
which is read ``if $\oe$ is of the form $\oinj{\ol_i}{\ox_i}$ then
$\op_i$ holds $(i = 1, \ldots, n)$'', with the understanding that the
expression is false if $\oe$ does not match any case. This can be
expressed with the negative formula
%
\begin{multline*}
  (\forall \ox_1.\, \oimply{\oe = \oinj{\ol_1}{\ox_1}}{\op_1})
  \land \cdots \land
  (\forall \ox_n.\,\oimply{\oe = \oinj{\ol_n}{\ox_n}}{\op_n})
  \land {} \\
  \lnot ((\forall x_1.\, \oe \neq \oinj{\ol_1}{\ox_1}) \land \cdots \land
  (\forall x_n.\, \oe \neq \oinj{\ol_n}{\ox_n})).
\end{multline*}
%
\fi %\iflong
%
The negative fragment is the part of first-order logic that has no
computational content in the realizability interpretation.
Consequently, the classical and constructive interpretations of
assertions agree. This is quite desirable, since RZ acts as a bridge
between constructive mathematics and real-world programmers, which
typically are not familiar with constructive logic.


RZ ever produces only a small subset of OCaml types (the unit type,
products, function types, polymorphic variant types, and parameter
types). Correspondingly, the language of terms produced is fairly simple
(tuples, functions, polymorphic variants, and local definitions).
However, the programmer is free to implement a specification using any
types and terms that exist in OCaml. 

A special kind of term is an
\emph{obligation} $\ooblig{\ox}{\oty}{\op}{\oe}$ which means ``in term
$e$, let $x$ be any element of $\values{\oty}$ that satisfies~$\op$''.
An obligation is equivalent to a combination of Hilbert's indefinite
description operator and a local definition,
$\olet{\ox}{(\varepsilon \ox {:} \oty.\,\op)}{\oe}$, where
$\varepsilon \ox {:} \oty.\, \op$ means ``any $x \in \values{\oty}$
such that $\op$''. The alternative form $\oobligx{\op}{\oe}$ stands
for $\ooblig{\_}{\ounit}{\op}{\oe}$. 

Obligations arise from the fact that well-formedness 
of the input language is undecidable; see Section~\ref{sec:input-language}.  
In such cases the system computes a realizability translation, but also produces obligations 
to be checked.  The programmer must
replace each obligation with a value satisfying the
obligation (i.e., demonstrate that the obligation can be satisfied). 
If such values do not exist, the specification is
unimplementable.

\else % \iflong

The output of RZ consists of \emph{module specifications}, module types
plus assertions about their components.  More specifically, a typical specification
may contain value declarations, type declarations and definitions, module
declarations, specification definitions, proposition declarations, and
assertions. 

The least familiar construct is the 
\emph{obligation} $\ooblig{\ox}{\oty}{\op}{\oe}$ which means ``in term
$e$, let $x$ be any element of $\values{\oty}$ that satisfies~$\op$''.
An obligation is equivalent to a combination of Hilbert's indefinite
description operator and a local definition,
$\olet{\ox}{(\varepsilon \ox {:} \oty.\,\op)}{\oe}$, where
$\varepsilon \ox {:} \oty.\, \op$ means ``any $x \in \values{\oty}$
such that $\op$''. The alternative form $\oobligx{\op}{\oe}$ stands
for $\ooblig{\_}{\ounit}{\op}{\oe}$. 

Obligations arise from the fact that well-formedness 
of the input language is undecidable; see Section~\ref{sec:input-language}.  
In such cases the system computes a realizability translation, but also produces obligations.  The programmer must
replace each obligation with a value satisfying the
obligation%
\iflong
(i.e., demonstrate that the obligation can be satisfied)
\fi % \iflong
. 
If such values do not exist, the specification is
unimplementable.

\fi % \iflong

%%% Local Variables: 
%%% mode: latex
%%% TeX-master: "cie"
%%% End: 

% Macros were moved to macros.tex because I need them in an earlier section.

\section{The Input Language}
\label{sec:input-language}

A RZ \emph{theory} is a generalized logical signature with associated
axioms, similar to a Coq module signature. Theories are descriptions
of \emph{models}, or implementations. The input is the definition of
one or more \emph{theories}; Figure~\ref{fig:kuratowski} shows a
sample input for the RZ system. RZ does not permit models to be
written directly, but it does allow a declaration of a model of a
given theory.

The simplest theory $\ith$ is a list of \emph{theory element}\/s
$\ite_1 \ldots \ite_n$. A theory element may specify that a certain
set, set element, proposition or predicate, or model must exist (i.e.,
\textsf{Parameter} specifications). It may also provide a definition
of a set, term, proposition, predicate, or theory (i.e.,
\textsf{Definition} specifications). Finally, a theory element can be
a named axiom that must hold. (i.e., \textsf{Axiom} specifications).

A theory may be parametrized by models of another theory,
$\ilambda{\iM}{\ith_1}{\ith_2}$. A typical example in mathematics is
the theory of a vector space, parameterized by a field of scalars. In
Figure~\ref{fig:kuratowski}, \texttt{K} is a theory parameterized by a
model choosing an implementation of a single set. For any such
model~$\im$, then, $\mathtt{K}(\im)$ is a theory describing a model of
finite sets whose elements are in the set $\iselect{\im}{\mathtt{a}}$.

A theory may also describe a uniform family of models
$\idarrow{\im}{\ith_1}{\ith_2}$, i.e., a single implementation of a
mapping\footnote{In OCaml such mappings are called \emph{functors}.}
that maps any model~$\im$ satisfying~$\ith_1$ to a model of~$\ith2$.
In contrast, a parameterized theory $T =
\ilambda{\iM}{\ith_1}{\ith_2}$ permits completely unrelated
implementations for $T(\im_1)$ and $T(\im_2)$. An example of a uniform
family of models would be a mapping from fields to vectors spaces that
maps a field~$F$ to the 1-dimensional vector space over~$F$.

\internal{Andrej}{I wanted to clarify the previous two paragraphs a
  bit, but I am not sure it's an improvement. Please improve some
  more. At some point we should say that both, parameterized theories
  and uniform families get translated to functor signatures. But where
  is the best place for that?}

\begin{figure}
	\[
	\begin{array}{rl@{\qquad}l}
		\noalign{\textbf{Propositions/Predicates}}
		\ip,\ipp ::= 
		    & \iP \ibar \iselect{\im}{\iP}&\mbox{Predicate/proposition names}\\
		  | & \itrue \ | \ \ifalse \ibar \inot{\ip_1} \ibar \iand{\ip_1}{\ip_2} \ibar 
		       \ior{\ip_1}{\ip_2}\ | \ \iimply{\ip_1}{\ip_2} & \mbox{Predicate logic}\\
		  | & \ilambda{\ix}{\is}{\ip}  \ibar \iapp{\ip}{\ie} & \mbox{Predicates and application}\\
		  | & \iequal{\ie_1}{\ie_2} & \mbox{Term equality}\\
		  | & \iforall{\ix}{\is}{\ip}  \ibar 
		      \iexists{\ix}{\is}{\ip} \ |\
		      \iunique{\ix}{\is}{\ip} & \mbox{Term quantifiers}\\[5pt]
		
		\noalign{\textbf{Sets}}
		\is ::= 
		    & \iS  \ibar \iselect{\im}{\iS} &\mbox{Set names}\\
		  | & \isunit \ibar \iprod{\ix}{\is_1}{\is_2}
                  &\mbox{Unit and (dependent) cartesian product}\\
		  | & \isvoid \ibar \isum{\il_1}{\is_1}{\il_2}{\is_2} &\mbox{Void and disjoint union}\\
		  | & \idarrow{\ix}{\is_1}{\is_2} & \mbox{(Dependent) function space} \\
		  | & \ilambda{\ix}{\is_1}{\is_2} \ibar 
		      \iapp{\is}{\ie} &\mbox{Dependent set and application}\\
		  | & \iquot{\is}{\ipp} & \mbox{Set quotient by an equivalence relation}\\
		  | & \isubset{\ix}{\is}{\ipp} & \mbox{Subset satisfying a predicate}\\
		  | & \irz{\is}&\mbox{Realizers of a set}\\[5pt] 
		
		\noalign{\textbf{Terms}}	
		\ie ::=
		    & \ix \ibar \iselect{\im}{\ix} &\mbox{Term names}\\
		  | & \ilambda{\ix}{\is_1}{\ie} \ibar 
		      \iapp{\ie_1}{\ie_2} &\mbox{Function and application}\\
		  | & \ituple{\ie_1}{\ie_2} 
		      \ibar \iproj{\ie}{n}&\mbox{Tuple and projection}\\
		  | & \iinj{\il}{\ie} 
		      \ibar (\imatch{\ie_0}{\il_1}{\ix_1}{\ie_1}{\il_2}{\ix_2}{\ie_2})&\mbox{Injection and projection from a union}\\
		  | & \ieclass{\ie}{\ipp}
		      \ibar \ileteclass{\ix}{\ipp}{\ie_1}{\ie_2}&\mbox{Equivalence class and picking a representative}\\
		  | & \irz{\ie}
		      \ibar \iletrz{\ix}{\ie_1}{\ie_2}&\mbox{Realized value and picking a realizer}\\
		  | & \isubin{\ie}{\is} \ibar \isubout{\ie}{\is} &\mbox{Type coercion in and out of a subset}\\
		  | & \ithe{\ix}{\is}{\ip}&\mbox{Definite description}\\
		  | & \ilet{\ix}{\ie_1}{\ie_2}&\mbox{Local definition}\\[5pt]
		
		\noalign{\textbf{Models}}		
		\im ::= 
		    & \iM  \ibar \iselect{\im}{\iM}&\mbox{Model names}\\
		  | & \iapp{\im_1}{\im_2}&\mbox{Application of parameterized model}\\[5pt]
		
		\noalign{\textbf{Proposition Kinds}}
		\ipt ::=
		    & \iProp & \mbox{Classifier for propositions}\\
		  | & \iStable & \mbox{Classifier for stable propositions}\\
		  | & \iEquiv{\is} &\mbox{Classifier for equivalence relations}\\
		  | & \idarrow{\ix}{\is}{\ipt} & \mbox{Classifier for a predicate/relation}\\[5pt] 
		
		\noalign{\textbf{Set Kinds}}
		\ik ::= 
		    & \iSet &\mbox{Classifier for a proper set}\\
		   | & \idarrow{\ix}{\is}{\ik} &\mbox{Classifier for a dependent set}\\[5pt]
		

		\noalign{\textbf{Theory Elements}}
		\ite ::=
		     & \iDefinition{\ix}{\ie}.&\mbox{Name for a term}\\
		   | & \iDefinition{\iS}{\is}.&\mbox{Name for a set}\\
		   | & \iDefinition{\iP}{\ip}.&\mbox{Name for a proposition/predicate}\\
		   | & \iDefinition{\iTH}{\ith}.&\mbox{Name for a theory}\\
		   | & \iParameter{\ix}{\is}.&\mbox{Require an element in the given set}\\
		   | & \iParameter{\iS}{\ik}.&\mbox{Require a set of the given kind}\\
		   | & \iParameter{\iP}{\ipt}.&\mbox{Require a proposition/predicate of the given sort}\\
		   | & \iParameter{\iM}{\ith}.&\mbox{Require a model of the given theory}\\
		   | & \iAxiom{\iP}{\ip}.&\mbox{Axiom that must hold}\\[5pt]

  		\noalign{\textbf{Theory}}
		\ith ::= 
		     & \iTH\\
%		   | & \iselect{\im}{\iTH}\\
		   	| & \ithy{\ite_1,\ldots,\ite_n}\\
		 	| & \idarrow{\iM}{\ith_1}{\ith_2} &\mbox{Theory of a uniform family of models}\\
		  	| & \ilambda{\iM}{\ith_1}{\ith_2} \ibar 
		      \iapp{\ith}{\iM}\\
	\end{array}
	\]
	\label{fig:input}
	\caption{Input Syntax (Simplified)}
\end{figure}

Propositions and predicates appearing in theories may use full
first-order constructive logic, not just the negative fragment. The
grammar for logical inputs is shown in Figure~\ref{fig:input}. Most of
this should be familiar, including the use of lambda abstraction to
define predicates.

The language of sets is rich, going well beyond the type systems of
typical programming languages. In addition to any base sets postulated
in a theory, one can construct dependent cartesian products and
dependent function spaces. We also supports disjoint unions (with
labeled tags), quotient spaces (a set modulo a stable equivalence
relation), subsets (elements of a set satisfying a predicate). RZ even
permits explicit references to sets of realizers.

The term language follows the structure of the set level by providing
introduction and elimination constructs. For product sets we have
tuples and projections ($\iproj{\ie}{1}$, $\iproj{\ie}{2}$, \ldots),
and for function spaces we have lambda abstractions and application.
One can inject a term into a tagged disjoin union, or do case analyses
on the members of a union. We can produce the equivalence class of a
given term, or pick a representative of a given equivalence class (as
long as what we do with it does not depend on the particular
representative). Very similarly, we can produce the set of realizers
for a term, or choose a representative from a given set of realizers
(as long as what we do with it does not depend on the particular
representative). We can inject a term into a subset (if it satisfies
the appropriate predicate), or project an element of a subset into the
enclosing set. Finally, the term language also allows local
definitions of term variables, and definite descriptions (as long as
there is a unique element satisfying the predicate in question).

From the parenthesized qualifiers in the previous paragraph, it is
clear that checking the well-formedness of terms is not decidable. RZ
checks what it can, but does not attempt any serious theorem proving.
Uncheckable constraints remain as explicitly-marked \emph{obligations}
in the final output, and must be verified by other means; the output
can then be trusted to the extent that the obligations are satisfied.

%%% Local Variables: 
%%% mode: latex
%%% TeX-master: "cie"
%%% End: 


\section{Translation}
\label{sec:translation}

Having shown the input and output languages for RZ,
%
\iflong
we now explain how the translation from one to the other works.
\else
we briefly sketch the translation from one to the other.
\fi
%
A theory is translated to a
specification, where the theory elements are translated as follows.



\subsection{Translation of sets and terms}
\label{sec:transl-sets-terms}

A set declaration $\iParameter{s}{\iSet}$ is translated to
%
\begin{source}
type s
predicate (\iPer{s}) : s \iTo s \iTo bool
assertion symmetric_s :  \iForall x:s, y:s, x \iPer{s} y -> y \iPer{s} x
assertion transitive_s : \iForall x:s, y:s, z:s, x \iPer{s} y \iAnd y \iPer{s} z \iTo x \iPer{s} z
predicate \iT{s} : s \iTo bool
assertion support_def_s :  \iForall x:s,  x : \iT{s} \iIff x \iPer{s} x
\end{source}
%
This says that the programmer should define a type~$s$ and a per
$\per_s$ on~$\values{s}$. Here~$\per_s$ is \emph{not} an OCaml value
of type $s \to s \to \mathtt{bool}$, but an abstract relation on the
set $\values{s} \times \values{s}$.  The relation may be uncomputable.

\iflong
The translation of the declaration of a dependent set
$\iParameter{\f{t}}{\f{s} \to \iSet}$ follows the interpretation of dependent
sets as uniform families (Section~\ref{sec:uniform-families}):
%
\begin{source}
type t
predicate \iPer{t} : s \iTo t \iTo t \iTo bool
assertion strict_t :  \iForall x:s, y:t, z:t, y \iPer{t x} z \iTo x : \iT{s}
assertion extensional_t :
  \iForall x:s, y:s, z:t, w:t, x \iPer{s} y \iTo z \iPer{t x} w \iTo z \iPer{t y} w
assertion symmetric_t : \iForall x:s, y:t, z:t, y \iPer{t x} z \iTo z \iPer{t x} y
assertion transitive_t :
  \iForall x:s, y:t, z:t, w:t, y \iPer{t x} z \iAnd z \iPer{t x} w \iTo y \iPer{t x} w
predicate \iT{t} : s \iTo t \iTo bool
assertion support_def_t :  \iForall x:s, y:t, y : \iT{t x} <\iTo y \iPer{t x} y
\end{source}
%
The underlying output type~$\f{t}$ is still non-dependent, but the per is
parameterized by~$\f{s}$.
\else % \iflong
The translation of the declaration of a dependent set
$\iParameter{\f{t}}{\f{s} \to \iSet}$ uses uniform families
(Section~\ref{sec:uniform-families}). The underlying type~$\f{t}$ is
non-dependent, but the per $\per_{\f{t}}$ receives an additional
parameter $\f{x} : \values{\f{s}}$.
\fi % \iflong

A value declaration $\iParameter{x}{s}$ is translated to
%
\begin{source}
val x : s
assertion x_support : x : \iT{s}
\end{source}
%
which requires the definition of a value~$\f{x}$ of type~$\f{s}$ which is in
the support of~$\f{s}$.
\iflong
When $\f{s}$ is not a basic set, RZ computes the
interpretation of the underlying type and support.
\fi

A value definition $\iDefinition{\f{x}}{e}$ where~$e$ is an expression
denoting an element of~$\f{s}$ is translated to
%
\newcommand{\iE}{$e$}
\begin{source}
val x : s
assertion x_def : x \iPer{s} \iE
\end{source}
%
The assertion does \emph{not} force~$\f{x}$ to be defined as~$e$, only
to be equivalent to it with respect to~$\per_{\f{s}}$. This is
useful, as often the easiest way to define a value is not the most
efficient way to compute it.

Constructions of sets in the input language are translated to
corresponding constructions of modest sets.
%
\iflong
In Section~\ref{sec:modest-sets-pers} we saw how products,
exponentials and their dependent versions are formed. We briefly
review the remaining constructions of modest sets. We only consider
those constructions of terms that are not entirely straightforward.
%
\else
%
We comment on those that are least familiar.
%
\fi

\iflong
\paragraph{Disjoint union.}
%
A disjoint union of modest sets $\isum{\il_1}{A}{\il_2}{B}$ is the
modest set whose underlying type is the sum of underlying types,
%
\begin{equation*}
  \typeOf{\isum{\il_1}{A}{\il_2}{B}} =
  \osumtyx{\il_1}{\typeOf{A}}{\il_2}{\typeOf{B}},
\end{equation*}
%
and the per is the disjoint union of pers $\per_{A}$ and $\per_{B}$, so
that we have
%
\begin{align*}
  \oper{\isum{\il_1}{A}{\il_2}{B}}{\oinj{\il_1}{u}}{\oinj{\il_1}{v}}
  &\iff \oper{s_1}{u}{v},
  \\
  \oper{\isum{\il_1}{A}{\il_2}{B}}{\oinj{\il_2}{u}}{\oinj{\il_2}{v}}
  &\iff \oper{s_2}{u}{v}.
\end{align*}
\fi % \iflong

\paragraph{Subsets.}
%
\iflong
%
The construction of subsets may look surprising at first, but
it makes sense computationally.
\fi
Given a predicate~$\phi$ on a per~$A$, the sub-per $\set{x : A
  \such \phi}$ has underlying type $\typeOf{A} \times \typeOf{\phi}$
where $(u_1,v_1) \per_{\set{x : A \such \phi}} (u_2,v_2)$ when $u_1
\per_A u_2$, $v_1 \rz \phi(u_1)$ and $v_2 \rz \phi(u_2)$. The point is
that a realizer for an element of $\set{x : A
  \such \phi}$ carries information about \emph{why} the element
belongs to the subset.

A type coercion $\icoerce{e}{t}$ can convert an element of the subset $s =
\set{x : t \such \phi(x)}$ to an element of~$t$. At the level of
realizers this is achieved by the first projection, which keeps a
realizer for the element but forgets the one for~$\phi(e)$. The
opposite type coercion $\icoerce{e'}{s}$ takes an $e' \in t$ and converts it
to an element of the subset. This is only well-formed when $\phi(e')$
is valid. Then, if $u \rz_t e'$ and $v \rz \phi(e')$, a realizer for
$\icoerce{e'}{s}$ is $(u, v)$. However, since RZ cannot in general know
a~$v$ which validates~$\phi(e')$, it emits the pair $(u,
(\ooblig{v}{\typeOf{\phi}}{\oapp{\oapp{\phi}{u}}{v}}{v}))$.

\paragraph{Quotients.}
%
\iflong
%
The category of modest sets has coequalizers, hence a quotient modest
set $A/\rho$ may be constructed for an any equivalence relation~$\rho$
on~$A$. However, because equality does not carry any computational
content, equality of equivalence classes $\ieclass{x}{\rho} =
\ieclass{y}{\rho}$ implies only $\inot{\inot{\rho(x,y)}}$, not the
usual~$\rho(x,y)$. As this may cause confusion
and mistakes, it is better to permit only quotients by \emph{stable}
equivalence relations, which behave as expected.

%
\else
%
Even though we may form quotients of pers by arbitrary equivalence
relations, only quotients by $\lnot\lnot$-stable relations behave as
expected.\footnote{The trouble is that from equality of equivalence
  classes $[x]_{\rho} = [y]_{\rho}$ we may conclude only
  $\lnot\lnot\rho(x,y)$ rather than the expected $\rho(x,y)$.}
%
\fi
A stable equivalence relation on a per~$A$ is the same thing as a
partial equivalence relation~$\rho$ on~$\typeOf{A}$ which satisfies
$\rho(x,y) \implies x \per_A y$. Then the quotient $A/\rho$ is the per
with $\typeOf{A/\rho} = \typeOf{A}$ and $x \per_{A/\rho} y \iff
\rho(x,y)$.

Luckily, it seems that many equivalence relations occurring in
computable mathematics are stable, or can be made stable with a little
bit of manipulation.
%
\iflong
For example, Cauchy sequences $(a_n)_{n\in\NN}$
and $(b_n)_{n\in\NN}$ represent the same real number when
%
\begin{equation}
  \label{eq:cauchy}%
  \xall{i}{\NN}{
    \xsome{j}{\NN}{
      \forall m,n \geq j.\,|a_m - b_n| \leq 2^{-i}
    }
  } \;.
\end{equation}
%
This defines an equivalence relation on the set of Cauchy sequences
which does not seem to be stable; intuitively a realizer for this
equivalence would be a computation telling us at what point
in the sequence the terms will be within~$2^{-i}$ of each other.
%
However, if we restrict attention just to the \emph{rapid} Cauchy
sequences, i.e., those satisfying $\xall{i}{\NN}{|a_{i+1} - a_i| \leq
  2^{-i}}$, then the equivalence relation becomes
%
\begin{equation*}
  \xall{i}{\NN}{|a_i - b_i| \leq 2^{-i+3}},
\end{equation*}
%
which is a negative formula; the above realizer is rendered
unnecessary.
%
\else % iflong
%
For example, the coincidence relation on Cauchy sequences is expressed
by a $\forall\exists\forall$~formula, but if we restrict to the
\emph{rapid} Cauchy sequences, it becomes a (negative)
$\forall$~formula.
%
\fi % iflong
It is interesting that most practical implementations of
real numbers follow this line of reasoning and represent real numbers in
way that avoids annotating every sequence with its rate of convergence.

Translation of an equivalence class $\ieclass{e}{\rho}$ is quite
simple, since a realizer for~$e$ also realizes its equivalence class
$\ieclass{e}{\rho}$. The elimination term
$\ileteclass{x}{\rho}{\xi}{e}$, means ``let $x$ be any element of
$\rho$-equivalence class $\xi$ in~$e$''. It is only well-formed when
$e$ does not depend on the choice of~$x$, but this is something RZ
cannot check. Therefore, if $u$ realizes $\xi$, RZ uses $u$ as a
realizer for~$x$ and emits an obligation saying that the choice of a
realizer for~$x$ does not affect~$e$.

\paragraph{The underlying set of realizers.}
%
Another construction on a per~$A$ is the underlying
per of realizers $\irz{A}$, defined by
%
\iflong
\begin{align*}
  \typeOf{\irz{A}} &= \typeOf{A} \\
  u \per_{\irz{A}} v &\iff u \in \support{A} \land u = v,
\end{align*}
%
\else  % \iflong
$\typeOf{\irz{A}} = \typeOf{A}$ and 
$u \per_{\irz{A}} v u \in \support{A} \land \iff u = v$, 
\fi % \iflong
%
where by $u = v$ we mean observational equality of values~$u$ and~$v$.
%
An element $r \in \irz{A}$ realizes a unique element $\irz{r} \in A$.
The elimination term $\iletrz{x}{e_1}{e_2}$, which means ``let $x$ be
any realizer for~$e_1$ in $e_2$'', is only well-formed if $e_2$ does
not depend on the choice of~$x$. This is an uncheckable condition,
hence RZ emits a suitable obligation in the output, and uses for $x$
the same realizer as for~$e_1$.

The construction $\irz{A}$ validates the Presentation Axiom (see
Section~\ref{sec:axiom-choice}). In the input language it gives us
access to realizers, which is useful because many constructions in
computable mathematics, such as those in Type Two
Effectivity~\cite{Wei00}, are explicitly expressed in terms of
realizers.

\iflong
\paragraph{Definite description.}
%
Russell's definite description operator $\ithe{x}{s}{\phi(x)}$ denotes
the unique element of $\isubset{x}{s}{\phi(x)}$. In case such an~$x$
does not exist, or if there are several, the term is not well formed.
The RZ translation essentially just asks the programmer to provide
suitable realizers for $x$ and for $\phi(x)$, and to check uniqueness,
%
\begin{equation*}
  \ooblig{x{:}s,b}{\typeOf{\phi}}{
    (\ototal{x}{s} \land \oapp{\oapp{\phi}{x}}{b} \land
    \oforall{x'}{s}{
      \oforall{c}{\typeOf{\phi}}{
        (\oimply{\oapp{\oapp{\phi}{x'}}{c}}{\oper{s}{x}{x'}})
      }
    })
    }{(x,b)}.
\end{equation*}
%
This is the best RZ can do, since in general it can check neither
that~$x$ exists, nor that it is unique.
\fi

\subsection{Translation of propositions}
\label{sec:transl-prop}

The driving force behind the translation of logic is a theorem
\cite[4.4.10]{Troelstra:van-Dalen:88:1} that says that under the
realizability interpretation every formula~$\phi$ is equivalent to one
that says, informally speaking, ``there exists~$u \in \typeOf{\phi}$,
such that $u$ realizes $\phi$''. Furthermore, the formula ``$u$
realizes $\phi$'' is computationally trivial. The translation of a
predicate $\phi$ then consists of its underlying type $\typeOf{\phi}$
and the relation $u \rz \phi$, expressed as a negative formula.

Thus an axiom $\iAxiom{\f{A}}{\phi}$ in the input is translated to
%
\newcommand{\iTypeOfPhi}{$|\phi|$}
\newcommand{\iPhi}{$\phi$}
\begin{source}
val r : \iTypeOfPhi
assertion A : u \iRz \iPhi
\end{source}
%
which requires the programmer to validate~$\phi$ by providing a
realizer for it. When $\phi$ is a compound statement RZ computes the
meaning \iflong of $u \rz \phi$ \fi as described in
Figure~\ref{fig:rz-logic}.

In RZ we avoid the explicit realizer notation $u \rz \phi$ in order to
make the output easier to read.
%
\iflong
%
A basic predicate declaration $\iParameter{p}{\iarrow{s}{\iProp}}$ is
translated to
%
\begin{source}
type ty_p
predicate p : s \iTo ty_p \iTo bool
assertion strict_p :  \iForall x:s, a:ty_p, p x a \iTo x : \iT{s}
assertion extensional_p :
   \iForall x:s, y:s, a:ty_p, x \iPer{s} y \iTo p x a \iTo p y a
\end{source}
%
We see that the predicate~$p$ has gained an additional argument of
type $\f{ty\_p}$ (which the programmer is supposed to define in an
implementation), and we write $\oapp{\oapp{p}{x}}{a}$ instead of $a
\rz \oapp{p}{x}$. The two assertions require that $p$ be strict and
extensional with respect to $\per_s$.
%
\else % iflong
%
A basic predicate declaration $\iParameter{p}{\iarrow{s}{\iProp}}$ is
translated to a type declaration $\otyspec{\f{ty\_p}}$ and a predicate
declaration $\opropspec{p}{s \to \f{ty\_p} \to \oProp}$ together with
assertions that~$p$ is strict and extensional.
%
\fi % iflong

Frequently we know that a predicate is stable, which can be taken into
account when computing its realizability interpretation. For this
purpose the input language has the subkind $\iStable$ of $\iProp$.
When RZ encounters a predicate which is declared to be stable, such as
$p : \iarrow{s}{\iStable}$, it does not generate a declaration of
$\f{ty\_p}$ and it does not give $p$ an extra argument.

Another special kind in RZ input language is the kind $\iEquiv{s}$ of
stable equivalence relations on a set~$s$. When an equivalence
relation is declared with $\iParameter{p}{\iEquiv{s}}$, RZ will output
1assertions stating that $p$ is strict, extensional, 
reflexive, symmetric and transitive.



%%% Local Variables: 
%%% mode: latex
%%% TeX-master: "cie"
%%% End: 

\section{Implementation}
\label{sec:implementation}

The RZ implementation consists of several sequential passes.

After the initial parsing, a \emph{type reconstruction} phase checks
that the input is well-typed (and checks for well-formedness to the
extent that it is easily decidable), and if successful produces an
annotated result with all variables explicitly tagged with types. The
type checking phase uses a system of dependent types, with limited
subtyping (implicit coercions) for sum types and subset types. The
details are fairly standard, so are omitted here. One non-obvious
consequence of the realizability translation, however, is that the
subset types $\isubset{\ix}{\iS}{\iand{\ipp_1(\ix)}{\ipp_2(\ix)}}$ and
$\isubset{\ix}{\iS}{\iand{\ipp_2(\ix)}{\ipp_1(\ix)}}$ are not
equal, but only isomorphic in general. An
explicit coercion is required to go from one type to the other,
because subset values are pairs containing realizers for
$\iand{\ipp_1(\ix)}{\ipp_2(\ix)}$ and
$\iand{\ipp_2(\ix)}{\ipp_1(\ix)}$, and these realizers have
potentially different types $|\ipp_1(\ix)|\mathtt{*}|\ipp_2(\ix)|$ and
$|\ipp_2(\ix)|\mathtt{*}|\ipp_1(\ix)|$ respectively.

Next the realizability translation is performed as described in
Section~\ref{sec:translation}, producing interface code. The
flexibility of the full input language (e.g., $n$-ary sum types and
dependent product types) makes the translation code fairly involved,
and so it is performed in a ``naive'' fashion whenever possible. The
immediate result of the translation is not easily readable.
 
\internal{Chris}{Uh oh...Kuratowski has been moved to a later section.
Should we move this section after
the examples?}

Thus, three more passes simplify the output before it is displayed to
the user. A \emph{thinning} pass removes all references to trivial
realizers produced by stable formulas. For example, direct translation
of the $\mathtt{free}$ axiom in the output for Kuratowski-finite sets
yields a value specification for $\mathtt{free}$ of type
%
\begin{equation*}
  (\mathtt{A.a} \to \mathtt{S.s}) \to 
  (\mathtt{fin} \to \mathtt{S.s}) * (\ounit * (\mathtt{A.a} \to
  \ounit) *
  (\mathtt{fin} \to \mathtt{fin} \to \ounit))
\end{equation*}
%
where $\ounit$ is the unit (terminal) type classifying the trivial
realizer. Thinning replaces this by the isomorphic type
%
\begin{equation*}
  (\mathtt{A.a} \to \mathtt{S.s}) \to \mathtt{fin} \to \mathtt{S.s}
\end{equation*}
%
and appropriately modifies references to $\mathtt{free}$ in the assertions to account for this change in type.

Next, an \emph{optimization} pass applies an ad-hoc collection of
basic logical and term simplifications in order to make the output more readable. 
Logical simplifications include applications of truth table rules
($\iand{\itrue}{\ip}$ becomes $\ip$), detection of syntactically
identical premises and conclusions
($\iimply{\ip_1}{\iand{\ip_1}{\ip_2}}$ becomes
$\iimply{\ip_1}{\ip_2}$), and optimization of other common patterns we have
seen arise
($\iforall{\ix}{\is}{\iimply{(\iequal{\ix}{\ie})}{\ipp(\ix)}}$ becomes
$\ipp(\ie)$). We do not attempt real theorem proving 
so some redundancy may remain, but in practice the optimization pass
can help significantly.

Finally, the user can specify whether two optional steps occur.
RZ can optionally performs a \emph{phase-splitting} pass~\cite{harper+:popl90}. 
This is an ML-specific transformation that replace certain
uses of parameterized modules (a heavyweight language construct) by
parameterized types and polymorphic values. The idea is that although
functors map modules containing types and terms to other modules containing types
and terms, constraints on the programming language ensure that output types
depend only on input types (and not input terms).  Thus, we can split each
functor into a mapping from input types to output types, and then a separate
(polymorphic) term mapping input types and terms to an output term.

For example (ignoring
assertions for simplicity) the entire module
\begin{source}
module Free : functor (S : Semilattice) ->
                    sig
                      val free : (A.a -> S.s) -> fin -> S.s
                    end	
\end{source}   
appearing in the output of the Kuratowski example can be replaced by the single polymorphic function
\begin{source}
val free : 's -> ('s -> 's -> 's) -> (A.a -> 's) -> fin -> 's	
\end{source}
which replaces the module parameter \texttt{S} by two extra term arguments term (corresponding to the module components \texttt{S.zero} and \texttt{S.join}) 
and a type argument \texttt{'s} for the type of lattice elements (corresponding to the module input \texttt{S.s}).

The other optional transformation is a \emph{hoisting} pass which lifts obligations in the output out to top-level positions.  This can make it easier to see exactly what one is obliged to provide.  When identical obligations appear in separate subterms of a term, hoisting can lift and merge these obligations, reducing redundancy.  However, moving obligations far from where they are used can make it harder to see why the obligation is required at all (and hence how one might satisfy the obligation), and so hoisting is turned off by default.

For example, in the following input (extracted from a larger description of an ordered field)
\begin{source}
Parameter s : Set.
Parameter zero : s.
Parameter inverse : {x : s | not (x = zero)} -> s.

Parameter lt : s -> s -> Stable.
Definition positive (x:s) := lt zero x.
Axiom lt_irr: forall x:s, not (lt x x).

Axiom order_inv: forall x:s, positive x -> positive (inverse x).
\end{source}
the axiom translates to the assertion:
\begin{source}
	(**  Assertion order_inv = 
          forall (x:||s||),  positive x ->
            positive (inv (assure (not (x =real= zero)) in x))
   *).
\end{source}
Here the system has noticed that for \Verb|inverse x| to make sense, we must know that
\Verb|x| is non-zero.   This requires non-trivial theorem proving and hence remains as 
an obligation for the user.  

We must prove \Verb|not (x =s= zero)| not for all \Verb|x|, but only under
the premises in force where the obligation occurs.  This is slightly clearer when hoisting moves the
obligation to the top level, after which it could be verified in the same way as all other assertions:
\begin{source}
   (**  Assertion order_inv = 
          assure (forall (x:||real||),  positive x -> not (x =real= zero))
            in forall (x:||real||),  positive x -> positive (inv x)
   *)	
\end{source}


%%% Local Variables: 
%%% mode: latex
%%% TeX-master: "cie"
%%% End: 

\section{Examples}
\label{sec:examples}

In this section we look at several examples which demonstrate various
points of RZ. Unfortunately, serious examples from computable
mathematics take too much space\footnote{The most basic structure
  in analysis (the real numbers) alone 
  requires several operations and a dozen or more axioms.} and will have to
be presented separately. The main theme is that constructively
reasonable axioms yield computationally reasonable operations.

\subsection{Decidable sets}
\label{sec:decidable-sets}

A set $S$ is said to be decidable when, for all $x, y \in S$, $x = y$
or $\lnot (x = y)$. In classical mathematics all sets are decidable, 
\iflong
because decidability of equality is just an instance of the law of
excluded middle.  But
\else
but
\fi % \iflong
RZ requires an axiom
%
\begin{source}
Parameter s : Set.
Axiom eq: \iForall x y : s, x = y \iOr \iNot (x = y).
\end{source}
%
to produce a realizer for equality
%
\begin{source}
val eq : s \iTo s \iTo [`or0 | `or1]
Assertion eq = \iForall (x:\iT{s}, y:\iT{s}),
                   (match eq x y with
                      `or0 \iImply x \iPer{s} y
                    | `or1 \iImply \iNot (x \iPer{s} y) )
\end{source}
%
We read this as follows: $\f{eq}$ is a function which takes
arguments~$\f{x}$ and~$\f{y}$ of type~$\f{s}$ and returns
$\mathtt{`or0}$ or $\mathtt{`or1}$. If it returns $\mathtt{`or0}$,
then $\oper{\f{s}}{\f{x}}{\f{y}}$, and if it returns
$\mathtt{`or1}$, then $\onot{(\oper{\f{s}}{\f{x}}{\f{y}})}$. In
other words $\f{eq}$ is a decision procedure%
\iflong
which tells when
values~$\f{x}$ and~$\f{y}$ represent the same element of the modest
set.
\else % \iflong
.
\fi % \iflong

\iflong
\subsection{Examples with obligations}
\label{sec:exampl-with-oblig}

In this section we show how RZ produces obligations, is sometimes able
to optimize them away, and show the effect of hoisting.

Consider how we might define division of real numbers. Assuming the
set of real numbers~$\f{real}$, a constant $\f{zero}$, and
multiplication operation~$\f{*}$ have already been declared and
axiomatized, we might write:
%
\begin{source}
Definition nonZeroReal := \{x : real | \iNot (x = zero)\}.
Parameter inv : nonZeroReal \iTo real.
Axiom inverse : \iForall x : real, \iNot (x = zero) -> x * (inv x) = one.
Definition (/) (x : real) (y : nonZeroReal) := x * (inv y).
\end{source}
%
We have defined the set of non-zero reals $\f{nonZeroReal}$ and
the inverse operation~$\f{inv}$ on it. Division $\f{x/y}$ is defined
as $\f{x * inv\; y}$. This does \emph{not} mean that the
programmer must necessarily implement division this way, only that the
implementation of $\f{x/y}$ must be equivalent to $\f{x * inv\;y}$.

In the axiom $\f{inverse}$, RZ encounters the subexpression
$\f{inv \;x}$. Because $\f{x}$ is quantified as an element of
$\f{real}$ rather than $\f{nonZeroReal}$, the typechecking
phase inserts a coercion that makes the expression well-typed.
Translation sees $\f{inv}(\f{x} \mathbin{{:}}
\f{nonZeroReal})$ instead of $\f{inv\ x}$ and translates this to
%
\begin{source}
inv (assure u:unit . \iNot (x \iPer{real} zero) in (x, u))
\end{source}
%
If this were the final output, the programmer would have to verify
that~$x$ is not zero, and provide a trivial realizer for it. However,
in this case the thinning phase first removes the trivial realizer,
%
\iflong
\begin{source}
inv (assure \iNot (x \iPer{real} zero) in x)
\end{source}%
\fi % \iflong
%
and then the optimizer determines that the obligation is not needed
because the whole expressions appears under the hypothesis that~$x$ is
not zero. So in the end the programmer sees
%
\begin{source}
(**  Assertion inverse =
  \iForall (x:\iT{real}),  \iNot (x \iPer{real} zero) \iTo (x * inv x) \iPer{real} one
*)
\end{source}
%
Assuming further that a strict linear order $<$ on~$\f{real}$ has
been axiomatized, we might proceed by relating it to $\f{inv}$:
%
\begin{showInput}
\iAxiom{\f{inv\_positive}}{
\iforall{\f{x}}{\f{real}}{\oimply{\f{zero} < \f{x}}{\f{zero} < \f{inv\ x}}}
}.
%Axiom inv_positive: forall x : real, zero < x -> zero < inv x.
\end{showInput}%
%
Once again $\f{inv\;x}$ appears in the input, bt this time the
optimizer is unable to remove the obligation, so the output is
%
\begin{showOutput}
\oassertion{\f{inv\_positive}}{}\\
\qquad
\oforallt{\f{x}}{\f{real}}{
  \oimply{
    \f{zero} < \f{x}}
  {\f{zero} < \f{inv}
    (\oobligx{\onot{(\oper{\f{real}}{\f{x}}{\f{zero}})}}{\f{x}})}
}
%(**  Assertion inv_positive =
%       forall (x:||real||),  zero < x ->
%         zero < inv (assure (not (x =real= zero)) in x)
%*)
\end{showOutput}
%
Local obligations can sometimes be hard to read, but if we activate the hoisting phase
(see Section~\ref{sec:implementation}), the obligation can be moved
to the top level. As this is done, the hypotheses under which the
obligation appears are collected, and we get
%
\begin{showOutput}
\oassertion{\f{inv\_positive}}{}\\
\qquad
\oobligx{
\oforallt{\f{x}}{\f{real}}{
\oimply{\f{zero}<\f{x}}{\onot{(\oper{\f{real}}{\f{x}}{\f{zero}})}}
}
}{}\\
\qquad\qquad{
\oforallt{\f{x}}{\f{real}}{
  \oimply{
    \f{zero} < \f{x}}
  {\f{zero} < \f{inv\ x}}
}}
%(**  Assertion inv_positive =
%  assure (forall (x:||real||),  zero < x -> not (x =real= zero))
%    in forall (x:||real||),  zero < x -> zero < inv x
%*)
\end{showOutput}%
%
Now it is easier to understand what must be checked, namely that
positive reals are not zero---an easy consequence of irreflexivity
of~$<$, but not something that RZ optimizer is aware of.

Lastly, we could define the golden ratio as the positive solution of
$x^2 = x + 1$,
%
\begin{source}
the x : real, (zero < x /\ x*x = x + one)
\end{source}
%
Not surprisingly, RZ cannot determine that there is a unique such~$x$,
so it outputs an obligation:
%
\begin{source}
assure x:real.
  (x : ||real|| /\ zero < x /\ x * x =real= x + one /\
     (forall (x':||real||),  zero < x' /\ x' * x' =real= x' + one ->
        x =real= x'))
  in x
\end{source}
\fi % \iflong

\iflong
\subsection{Finite sets}
\label{sec:finite-sets}

\begin{figure}[t]
\begin{showInputSmall}
\iDefinition{\f{Semilattice}}{\mathsf{thy}}\\
\qquad \iParameter{\f{s}}{\iSet}.\\
\qquad \iParameter{\f{zero}}{\f{s}}.\\
\qquad \iParameter{\f{join}}{\f{s}\to\f{s}\to\f{s}}.\\
\\
\qquad \iAxiom{\f{commutative}}{\forall \f{x},\f{y}:\f{s}.\ \f{join\ x\ y} = \f{join\ y\ x}}.\\
\qquad \iAxiom{\f{associative}}{\forall \f{x},\f{y},\f{z}:\f{s}.\ \f{join}\,(\f{join\ x\ y})\,\f{z} = \f{join}\,\f{x}\, (\f{join\ y\ z})}.\\
\qquad \iAxiom{\f{idempotent}}{\forall \f{x}:\f{s}.\ \f{join\ x\ x} = \f{x}}.\\
\qquad \iAxiom{\f{neutral}}{\forall \f{x}:\f{s}.\ \f{join\ x\ zero} = \f{x}}.\\
\mathsf{end}.
\end{showInputSmall}
  \caption{The theory of a semilattice}
  \label{fig:semilattice}
\end{figure}

\iflong
There are many characterizations of finite sets, but the one that
works best constructively is due to Kuratowski, who identified the
finite subsets of~$A$ as the least family~$K(A)$ of subsets of~$A$
that contains the empty set and is closed under unions with
singletons. This characterization relies on powersets, which are not
available in RZ. But the gist of it, namely that $K(A)$ is an
\emph{inital} structure a suitable sort, can be expressed as follows.

\else
%
The family $K(A)$ of finite subsets of a set~$A$ may be characterized
as the free $\vee$-semilattice generated by~$A$.
%
\fi
%
Recall that a \emph{$\vee$-semilattice} is a set~$S$ with a
constant~$0 \in S$ and an associative, commutative, and idempotent
operation ``join'' $\vee$ on~$S$ such that $0$ is the neutral element
for~$\vee$, see Figure~\ref{fig:semilattice} for RZ axiomatization of
semilattices.
%
The Kuratowski finite sets~$K(A)$ are the \emph{free} semilattice
generated by a set~$A$, where $\vee$ is union and $0$ is the empty
set. This is formalized in RZ as shown in Figure~\ref{fig:kuratowski}.
%
\begin{figure}
\begin{showInput}
\iDefinition{\texttt{K}\ (\texttt{A} : \f{thy}\ \  \iParameter{\f{a}}{\iSet}.\ \  \texttt{end})}{\f{thy}}\\
\qquad \iInclude{\f{Semilattice}}.\\
\qquad \iParameter{\f{singleton}}{\f{A.a} \to \f{s}}.\\
\qquad \iDefinition{\f{fin}}{\f{s}}.\\
\qquad \iDefinition{\f{emptyset}}{\f{zero}}.\\
\qquad \iDefinition{\f{union}}{\f{join}}.\\
\\
\qquad \iAxiom{\f{free}}{} \forall \f{S}:\f{Semilattice}.\ \forall\f{f}:{\f{A.a}\to\f{S.s}}.\ \exists!g:\f{fin}{\to}{\f{S.s}}.\\
\qquad \qquad \qquad \qquad \qquad \f{g\ emptyset} = \f{S.zero}\ \land\\
\qquad \qquad \qquad \qquad \qquad \forall\f{x}:\f{A.a}.\ \f{f\ x} = \f{g}\ (\f{singleton\ x})\ \land\\
\qquad \qquad \qquad \qquad \qquad \forall\f{u},\f{v}:\f{fin}.\ \f{g}(\f{union\ u\ v}) = \f{S.join}\,(\f{g\ u})\,(\f{g\ v}).\\
\f{end}.
\end{showInput}
  \caption{Kuratowski finite sets}
  \label{fig:kuratowski}
\end{figure}
%
The theory $\f{K}$ is parametrized by a model~$\f{A}$ which contains a
set~$\f{a}$. In the first line we include the theory of semilattices.
Then we postulate an operation $\f{singleton}$ which injects the
generators into the semilattice. The three definitions are just a
convenience, so that we can refer to the parts of $\f{K(A)}$ by their
natural names, e.g., $\f{emptyset}$ instead of $\f{zero}$. The axiom
$\f{free}$ expresses the fact that $\f{K}(\f{A})$ is the free
semilattice on~$\f{A.a}$: for every semilattice $\f{S}$ and a map
$\f{f} : \f{A.a} \to \f{S.s}$ from the generators to the underlying
set of~$\f{S}$, there exists a unique semilattice homomorphism $\f{g}
: \f{fin} \to \f{S.s}$ such that $\f{f}(\f{x}) = \f{g}(\f{singleton\;
  x})$.

The output for $\f{Semilattice}$ and~$\f{K}$ specifies
values of suitable types for each declared constant and operation. All
axioms but the last one are equations and have straightforward
translations in terms of underlying pers. The output for the axiom
$\f{free}$ is shown in Figure~\ref{fig:free}.
%
\begin{figure}
  \centering
\begin{showOutputSmall}
\f{module}\ \f{Free} : \f{functor} (\f{S} : \f{Semilattice}) \to \f{sig}\\
\quad \ovalspec{\f{free}}{\oarrow{(\oarrow{\f{A.a}}{\f{S.s}})}{\oarrow{\f{fin}}{\f{S.s}}}}
\\
\quad
\oassertion{\f{free}}{}
\\
\quad\quad
\oforallt{\f{f}}{\oarrow{\f{A.a}}{\f{S.s}}}{
\olet{\f{g}}{\f{free\ f}}{}}\\
\qquad\qquad \ototal{\f{g}}{\oarrow{\f{fin}}{\f{S.s}}} \land \oper{\f{S.s}}{\f{g\ emptyset}}{\f{S.zero}} \land {}\\
\qquad\qquad (\oforallt{\f{x}}{\f{A.a}}{\oper{\f{S.s}}{\f{f\ x}}{\f{g(singleton\ x)}}}) \land {} \\
\qquad\qquad (\oforallt{\f{u},\f{v}}{\f{fin}}{\oper{\f{S.s}}{\f{g(union\ u\ v)}}{\f{S.join\ (g\ u)\ (g\ v)}}}) \land {} \\
\qquad\qquad
((\oforall{\f{h}}{\oarrow{\f{fin}}{\f{S.s}}}{}\\
\qquad\qquad\qquad
\ototal{\f{h}}{\oarrow{\f{fin}}{\f{S.s}}} \land
\oper{\f{S.s}}{\f{h\ emptyset}}{\f{S.zero}} \land {} \\
\qquad\qquad\qquad (\oforallt{\f{x}}{\f{A.a}}{\oper{\f{S.s}}{\f{f\ x}}{\f{h(singleton\ x)}}}) \land {} \\
\qquad\qquad\qquad
(\oforallt{\f{u},\f{v}}{\f{fin}}{\oper{\f{S.s}}{\f{h(union\ u\ v)}}{\f{S.join\ (h\ u)\ (h\ v)}}})) \Rightarrow {} \\
\qquad\qquad
\oforall{\f{x},\f{y}}{\f{fin}}{\oimply{\oper{\f{fin}}{\f{x}}{\f{y}}}{\oper{\f{S.s}}{\f{g\ x}}{\f{h\ y}}}})
%module Free : functor (S : Semilattice) ->
%sig
%val free : (A.a -> S.s) -> fin -> S.s
%(**  Assertion free = 
%forall (f:||A.a -> S.s||), 
%  let g = free f in g : ||fin -> S.s|| /\ 
%  g emptyset =S.s= S.zero /\ 
%  (forall (x:||A.a||),  f x =S.s= g (singleton x)) /\ 
%  (forall (u:||fin||, v:||fin||), g (union u v) =S.s= S.join (g u) (g v)) /\ 
%  (forall h:fin -> S.s,  h : ||fin -> S.s|| /\ 
%     h emptyset =S.s= S.zero /\ 
%     (forall (x:||A.a||), f x =S.s= h (singleton x)) /\ 
%     (forall (u:||fin||, v:||fin||), 
%        h (union u v) =S.s= S.join (h u) (h v)) ->
%     forall x:fin, y:fin,  x =fin= y -> g x =S.s= h y)
%*)
%end
\end{showOutputSmall}%
  \caption{Output of axiom $\texttt{free}$.}
  \label{fig:free}
\end{figure}
%
Because the axiom quantifies over all models~$\f{S}$ of the theory
$\f{Semilattice}$ its translation is a functor~$\f{Free}$ which
accepts an implementation of a semilattice~$S$ and yields a realizer
$\f{free}$ validating the axiom. The computational meaning of
$\f{free}$ is a combination map and fold operation, taking a map
$\f{f} : \f{A.a} \to \f{S.s}$ and a finite set~$\f{u} = \set{x_1,
  \ldots, x_n}$, and return $\f{f}(x_1) \vee \cdots \vee \f{f}(x_n)$,
where $\vee$ is the join operation on the semilattice~$S$.

Applying phase-splitting to this axiom yields the even simpler
specification
%
\begin{equation*}
\mathtt{val}\ \f{free} : \alpha \to (\alpha \to \alpha \to \alpha) \to (\f{A.a}\to\alpha) \to \f{fin} \to \alpha	
\end{equation*}
%
(with an appropriate assertion)
which replaces the module parameter \texttt{S} by two extra term arguments term (corresponding to the module components \texttt{S.zero} and \texttt{S.join}) 
and a type argument $\alpha$ for the type of lattice elements (corresponding to the module input \texttt{S.s}).  This is even
more recognizable as a folding operation over the set.


It is important to note that, in contrast to \texttt{fold} operators found in typical functional
languages, \texttt{free} is only expected to work for suitable \texttt{join} arguments (e.g., idempotent and order independent).  These
sets are not the typical finite-set data structure: there is no membership predicate, nor
is there a way to compute the size of a set.  There is no
assumption that equality is decidable for set elements; this permits
finite sets of  exact real numbers, for example.  Decidable equality
is required both for membership and for detecting
whether the same element has been added twice to the same set\footnote{The natural implementation would thus
be an unordered collection of elements, possibly with duplicates.}.

Some operations are nevertheless computable.  Using \texttt{free} one
can determine whether a finite set is empty.  In the case of a set of exact
real numbers, we cannot compute their sum, but we could compute maximum or minimum.

More common set implementations (e.g., the \texttt{Set} module in the OCaml standard library)
implement sets over values with either decidable total order; these could also be
formalized in RZ.
\fi % \iflong

\subsection{Inductive types}
\label{sec:inductive-types}

To demonstrate the use of dependent types we show how RZ handles
general inductive types, also known as W-types or general
trees~\cite{nordstroem90:_progr_martin_type_theor}. Recall that a
W-type is a set of well-founded trees, where the branching types of
trees are described by a family of sets $B = \set{T(x)}_{x \in S}$.
Each node in a tree has a \emph{branching type}~$x \in S$, which
determines that the successors of the node are labeled by the elements
of~$T(x)$.
%
\iflong
%
For example, to get non-empty binary trees whose leaves are
labeled by natural numbers, define
%
\begin{align*}
  S &= \set{\f{cons}} \cup \set{\f{leaf}(n) \such n \in \NN}
  \\
  T(\f{cons}) &= \set{\f{left}, \f{right}}
  \\
  T(\f{leaf}(n)) &= \emptyset.
\end{align*}
%
Then a node of type $\f{cons}$ has two successors, indexed by
constants $\f{left}$ and $\f{right}$, while a node of type
$\f{leaf}(n)$ does not have any successors.
\par
%
\fi % iflong
%
Figure~\ref{fig:wtype} shows an RZ axiomatization of W-types.
%
\begin{figure}
\begin{source}
Parameter W : [B : Branching] \iTo
thy
  Parameter w : Set.
  Parameter tree : [x : B.s] \iTo (B.t x \iTo w) \iTo w.
  Axiom induction:
    \iForall M : thy Parameter p : w \iTo Prop. end,
    (\iForall x : B.s, \iForall f : B.t x \iTo w,
       ((\iForall y : B.t x, M.p (f y)) \iTo M.p (tree x f))) \iTo
    \iForall t : w, M.p t.
end.
\end{source}
  \caption{General inductive types}
  \label{fig:wtype}
\end{figure}
%
The theory $\f{Branching}$ describes that a branching type
consists of a set~$\f{s}$ and a set~$\f{t}$ depending on~$\f{s}$. The theory~$\f{W}$ is
parameterized by a branching type~$\f{B}$. It specifies a set~$\f{w}$ of
well-founded trees and a tree-forming operation $\f{tree}$ with a
dependent type $\Pi_{\f{x} \in \f{B.s}} (\f{B.t(x)} \to \f{w}) \to \f{w}$.
%
\iflong
%
Given a
branching type~$\f{x}$ and a map $\f{f} : \f{B.t(x)} \to \f{w}$, $\f{tree}\;\f{x}\;\f{f}$
is the tree whose root has branching type~$\f{x}$ and whose successor
labeled by $\ell \in \f{B.t}(\f{x})$ is the tree~$\f{f}(\ell)$.
%
\fi
%
The inductive nature of~$\f{w}$ is expressed with the axiom
$\f{induction}$, which states that for every property $\f{M.p}$, if $\f{M.p}$
is an inductive property then every tree satisfies it. A property is
said to be \emph{inductive} if a tree $\f{tree}\;\f{x}\;\f{f}$ satisfies it
whenever all its successors satisfy it.

\iflong
In the translation, see Appendix~\ref{sec:outp-induct-types} for a
complete output, dependencies at the level of types and terms disappear.
\else
In the translation dependencies at the level of types and terms disappear.
\fi
%
A branching type is determined by a pair of non-dependent types~$\f{s}$
and~$\f{t}$ but the per $\per_{\f{t}}$ depends on~$\values{\f{s}}$. The theory~$\f{W}$
turns into a signature for a functor receiving a branching type~$\f{B}$
and returning a type~$\f{w}$, and an operation $\f{tree}$ of type
$\f{B.s} \to (\f{B.t} \to \f{w}) \to \f{w}$.  One can use phase-splitting
to translate axiom
$\f{induction}$ into a specification of a polymorphic function
%
\begin{equation*}
  \f{induction:
  (B.s \to (B.t \to w) \to (B.t \to \poly{\alpha}) \to \poly{\alpha}) \to w \to \poly{\alpha}},
\end{equation*}
%
which is a form of recursion on well-founded trees. Instead of trying
to explain what $\f{induction}$ is supposed to do, we show a surprisingly simple,
hand-written implementation of W-types. The reader may enjoy figuring out how it works:
%
\sourcefile{wtype.ml}


\subsection{Axiom of choice}
\label{sec:axiom-choice}

RZ can help explain why a generally
accepted axiom is not constructively valid. Consider the Axiom of
Choice:
%
\begin{source}
Parameter a b : Set.
Parameter r : a \iTo b \iTo Prop.
Axiom ac: (\iForall x : a, \iExists y : b, r x y) \iTo
          (\iExists c : a \iTo b, \iForall x : a, r x (c x)).
\end{source}
%
The relevant part of the output is
%
\begin{source}
val ac : (a \iTo b * ty_r) \iTo (a \iTo b) * (a \iTo ty_r)
Assertion ac =
  \iForall f:a \iTo b * ty_r,
    (\iForall (x:\iT{a}),  let (p,q) = f x in p : \iT{b} \iAnd r x p q) \iTo
    let (g,h) = ac f in
      g : \iT{a \iTo b} \iAnd (\iForall (x:\iT{a}),  r x (g x) (h x))
\end{source}
%
This requires a function $\f{ac}$ which accepts a function $\f{f}$
and computes a pair of functions $\f{(g,h)}$. The input function~$\f{f}$ takes
an $\ototal{\f{x}}{\f{a}}$ and returns a pair $\f{(p,q)}$ such that $\f{q}$ realizes
the fact that $\f{r\;x\;p}$ holds. The output functions $\f{g}$ and $\f{h}$ taking
$\ototal{\f{x}}{\f{a}}$ as input must be such that $\f{h\;x}$ realizes
$\f{r\;x\;(g\;x)}$. Crucially, the requirement $\ototal{\f{g}}{\oarrow{\f{a}}{\f{b}}}$
says tht $\f{g}$ must be extensional, i.e., map equivalent realizers to
equivalent realizers. We could define~$\f{h}$ as the first component
of~$\f{f}$, but we cannot hope to implement~$\f{g}$ in general because the
second component of~$f$ is not assumed to be extensional.

The \emph{Intensional} Axiom of Choice allows the choice function to
depend on the realizers:
%
\begin{source}
Parameter a b : Set.
Parameter r : a \iTo b \iTo Prop.
Axiom iac: (\iForall x : a, \iExists y : b, r x y) \iTo
           (\iExists c : rz a \iTo b, \iForall x : rz a, r (rz x) (c x)).
\end{source}
%
Now the output is
%
\begin{source}
val iac : (a \iTo b * ty_r) \iTo (a \iTo b) * (a \iTo ty_r)
Assertion iac =
  \iForall f:a \iTo b * ty_r,
    (\iForall (x:\iT{a}),  let (p,q) = f x in p : \iT{b} \iAnd r x p q) \iTo
    let (g,h) = iac f in
      (\iForall x:a, x : \iT{a} \iTo g x : \iT{b}) \iAnd (\iForall (x:\iT{a}),  r x (g x) (h x))
\end{source}
%
which is exactly the same as before, \emph{except} that the
troublesome requirement $\ototal{\f{g}}{\oarrow{\f{a}}{\f{b}}}$ turned into
$\oforall{\f{x}}{\f{a}}{(\oimply{\ototal{\f{x}}{\f{a}}}{\ototal{\f{g\;x}}{\f{b}}})}$, which
is weaker. We can impement $\f{iac}$ as
%
\begin{source}
let iac f = (fun x -> fst (f x)), (fun x -> snd (f x))
\end{source}
%
The Intensional Axiom of Choice is in fact just an instance of the
usual Axiom of Choice applied to~$\irz{A}$ and~$B$. Combined with the
fact that~$\irz{A}$ covers~$A$, this establishes the validity of
\emph{Presentation Axiom}~\cite{barwise75:_admis_sets_struc}, which
states that every set is an image of one satisfying the axiom of
choice.

\subsection{Modulus of Continuity}
\label{sec:we-show-modulus-of-continuity-example}

As a last example we show how certain constructive principles require
the use of computational effects. To keep the example short, we
presume that we are already given the set of natural
numbers~$\f{nat}$ with the usual structure.

A \emph{type 2 functional} is a map $f : (\f{nat} \to \f{nat})
\to \f{nat}$. It is said to be continuous if the output of $f(a)$
depends only on an initial segment of the sequence~$a$. We can express
this axiom in RZ as follows:
%
\begin{source}
Axiom continuity: \iForall f : (nat \iTo nat) \iTo nat, \iForall a : nat \iTo nat,
  \iExists k, \iForall b : nat \iTo nat, (\iForall m, m \iLeq k \iTo a m = b m) \iTo f a = f b.
\end{source}
%
The axiom says that for any $\f{f}$ and $\f{a}$ there exists $\f{k} \in
\f{nat}$ such that $\f{f(b) = f(a)}$ as soon as the sequences~$\f{a}$
and~$\f{b}$ agree on the first $\f{k}$ terms. The axiom is translated to the
specification
%
\begin{source}
val continuity : ((nat \iTo nat) \iTo nat) \iTo (nat \iTo nat) \iTo nat
Assertion continuity =
  \iForall (f:\iT{(nat \iTo nat) \iTo nat}, a:\iT{nat \iTo nat}),
    let p = continuity f a in p : \iT{nat} \iAnd
    (\iForall (b:\iT{nat \iTo nat}),
       (\iForall (m:\iT{nat}),  m \iLeq p \iTo a m \iPer{nat} b m) \iTo f a \iPer{nat} f b)
\end{source}
%
which says that $\f{continuity\;f\;a}$ is a number~$\f{p}$ such that
$\f{f(a) = f(b)}$ whenever $\f{a}$ and $\f{b}$ agree on the first~$\f{p}$ terms. In
other words, $\f{continuity}$ is a \emph{modulus of continuity}
functional. It cannot be implemented in a purely functional
language,\footnote{There are models of $\lambda$-calculus which validate
  the choice principle~$AC_{2,0}$, but this principle contradicts the
  existence of a modulus of continuity functional,
  see~\cite[9.6.10]{Troelstra:van-Dalen:88:2}.} but with the use of
store we can implement it as
%
\begin{source}
let continuity f a =
  let p = ref 0 in
  let a' n = (p := max !p n; a n) in
    f a' ; !p
\end{source}
%
To compute a modulus for~$\f{f}$ at~$\f{a}$, the program creates a
function~$\f{a'}$ which is just like~$\f{a}$ except that it stores in~$\f{p}$ the
largest argument at which it has been called. Then $\f{f\;a'}$ is
computed, its value it discarded, and the value of~$\f{p}$ is returned.
The program works because~$\f{f}$ is assumed to be extensional and must
therefore not distinguish between extensionally equal sequences~$\f{a}$
and~$\f{a'}$.



%%% Local Variables: 
%%% mode: latex
%%% TeX-master: "cie"
%%% End: 

\section{Related Work}
\label{sec:related-work}

\subsection{Coq}
\label{sec:comparison-with-coq}

Coq is a very flexible system, with complete support for
theorem-proving and creating trusted code. One common pattern of use
for Coq is to write code in Coq's functional language (values whose
types are \texttt{Set} in Coq), to state and prove theorems stating
that the code behaves correctly (where the theorems are Coq values
whose types are \texttt{Prop} in Coq), and then have Coq produce
guaranteed correct code in ML. In such cases, RZ is complementary to
Coq. RZ can clarify the constructive content of mathematical
structures and hence suggest an appropriate division between Coq's
\texttt{Set} and \texttt{Prop}, i.e., which values and which theorems
should appear in the input to Coq. (It should be easy to have RZ
produce output in Coq syntax, and we hope to do this eventually.)

In general, RZ is a smaller and more lightweight system and thus more
flexible where it applies. It is not always practical or necessary to
do theorem proving in order to provide an implementation; interfaces
generated by RZ can be implemented in any manner from theorem proving
to directly writing code. And, RZ provides a way to talk with
programmers about implications for constructive without bringing in
full theorem proving.

\internal{Andrej}{I don't parse the previous sentence.}

\subsection{Other tools}

Komagata and Schmidt~\cite{komagata+:tr95} describe a system that uses
a similar realizability translation to ours. Like Coq, the system
translates formal proofs to programs, rather than

An interesting technical difference is that the algorithm they use,
attributed to John Hatcliff, does thinning as it goes along, rather
than making this a separate pass. For example, the translation of the
conjunction-introduction rule has four cases, depending on whether the
left and/or right propositions being proved are [almost?] negative, in
which case the trivial contribution can be immediately discarded.

\subsection{Other Models of Computability}
\label{sec:models-of-computability}

Most formulations of computable mathematics are based on realizability
models~\cite{Bauer:00}, even though they were not initially developed,
nor are they usually presented within the framework of realizability:
Recursive Mathematics~\cite{ershov98:_handb_recur_mathem} is based on
the original realizability by Turing machines~\cite{KleeneSC:intint};
Type Two Effectivity~\cite{Wei00} on function
realizability~\cite{KleeneSC:fouim} and relative function
realizability~\cite{BirkedalL:devttc}, while topological and domain
representations~\cite{Bla97a,Bauer:Birkedal:Scott:98} are based on
realizability over the graph model
$\mathcal{P}\omega$~\cite{ScottD:dattl}. A common feature of these is
that they use models of computation which are suitable for the
theoretical studies of computability, rather than for practical
programming. Since we want to emphasize the practical aspects of
realizability, we have chosen instead an actual real-world programming
language.


%%% Local Variables: 
%%% mode: latex
%%% TeX-master: "cie"
%%% End: 


%% Bibliography
% AB: It's weird that references are ordered by first appearance, but if
% that's what they want, that's what they get.
\bibliographystyle{splncs}
\bibliography{rzbib}

\appendix

\section{Inductive types}
\label{sec:outp-induct-types}

%\subsection{Inductive types}
\label{sec:inductive-types}

To demonstrate the use of dependent types we show how RZ handles
general inductive types, also known as W-types or general
trees~\cite{nordstroem90:_progr_martin_type_theor}. Recall that a
W-type is a set of well-founded trees, where the branching types of
trees are described by a family of sets $B = \set{T(x)}_{x \in S}$.
Each node in a tree has a \emph{branching type}~$x \in S$, which
determines that the successors of the node are labeled by the elements
of~$T(x)$. For example, to get non-empty binary trees whose leaves are
labeled by natural numbers, define
%
\begin{align*}
  S &= \set{\mathtt{cons}} \cup \set{\mathtt{leaf}(n) \such n \in \NN}
  \\
  T(\mathtt{cons}) &= \set{\mathtt{left}, \mathtt{right}}
  \\
  T(\mathtt{leaf}(n)) &= \emptyset.
\end{align*}
%
Then a node of type $\mathtt{cons}$ has two successors, indexed by
constants $\mathtt{left}$ and $\mathtt{right}$, while a node of type
$\mathtt{leaf}(n)$ does not have any successors.

Figure~\ref{fig:wtype} shows an RZ axiomatization of W-types.
%
\begin{figure}
  \centering
  \sourcefile{wtype.thy}
  \caption{General inductive types}
  \label{fig:wtype}
\end{figure}
%
The theory $\mathtt{Branching}$ describes that a branching type
consits of a set~$s$ and a set~$t$ depending on~$s$. The theory~$W$ is
parametrized by a branching type~$B$. It specifies a set~$w$ of
well-founded trees and a tree-forming operation $\mathtt{tree}$ with a
dependent type $\prod_{x \in B.s} (B.t(x) \to w) \to w$. Given a
branching type~$x$ and a map $f : B.t(x) \to w$, $\mathtt{tree}\,x\,f$
is the tree whose root has branching type~$x$ and whose successor
labeled by $\ell \in B.t(x)$ is the tree~$f(\ell)$. The inductive
nature of~$w$ is expressed with the axiom $\mathtt{induction}$, which
states that for every property $M.p$, if $M.p$ is an inductive
property then every tree satsifies it. A property is said to be
\emph{inductive} if a tree $\mathtt{tree}\,x\,f$ satisfies it whenever
all its successors satisfy it.

In the translation, see Appendix~\ref{sec:outp-induct-types},
dependencies at the level of types and terms disappear. A branching
type is determined by a pair of non-dependent types~$s$ and~$t$ but
the per $\per_t$ depends on~$\values{s}$. The theory~$W$ turns into a
signature for a functor receiving a branching type~$B$ and returning a
type~$w$, and an operation $\mathtt{tree}$ of type $B.s \to (B.t \to
w) \to w$. In this example we use phase-splitting (see
Section~\ref{sec:implementation}) to translate axiom
$\mathtt{induction}$ into a specification of a polymorphic function
(compare with output for axiom \texttt{free} in the previous example)
%
\begin{equation*}
  \mathtt{induction}:
  (B.s \to (B.t \to w) \to (B.t \to \poly{ty\_p}) \to \poly{ty\_p}) \to w \to \poly{ty\_p},
\end{equation*}
%
which is a form of recursion on well-founded trees. Instead of trying
to explain what $\mathtt{induction}$ is supposed to do, we show in
Figure~\ref{fig:wtype-implementation} a surprisingly simple, complete
hand-written implementation of W-types. The reader may be entertained
by figuring out how $\mathtt{induction}$ works.
%
\begin{figure}
  \centering
  \sourcefile{wtype.ml}
  \caption{An implementation of general inductive types.}
  \label{fig:wtype-implementation}
\end{figure}

The following is the unabridged output for the theory of inductive
types shown in Figure~\ref{fig:wtype}.

\sourcefile{wtype.mli}

%%% Local Variables: 
%%% mode: latex
%%% TeX-master: "cie.tex"
%%% End: 



\end{document}

%%% Local Variables: 
%%% mode: latex
%%% TeX-master: t
%%% End: 
